\documentclass[11pt]{article}
\usepackage[a4paper,margin=25mm]{geometry}
\usepackage{amsmath,amssymb,amsthm,booktabs,longtable,array,hyperref,microtype}
\hypersetup{colorlinks=true,linkcolor=blue,urlcolor=blue,pdftitle={Acquiring local algebras for exact counting},pdfauthor={Astra P3}}
\newtheorem{theorem}{Theorem}[section]\newtheorem{lemma}[theorem]{Lemma}
\newtheorem{proposition}[theorem]{Proposition}\newtheorem{corollary}[theorem]{Corollary}
\theoremstyle{remark}\newtheorem{remark}[theorem]{Remark}
\newcommand{\F}{\mathbb F}\newcommand{\Tr}{\operatorname{Tr}}
\newcommand{\rank}{\operatorname{rank}}\newcommand{\im}{\operatorname{im}}
\newcommand{\Cent}{\operatorname{Cent}}\newcommand{\poly}{\operatorname{poly}}
\newcommand{\Span}{\operatorname{span}}\newcommand{\GL}{\operatorname{GL}}
\title{Acquiring local algebras for exact counting\\
\large Quadratic weights, nilpotent depth, and critical-ray reconstruction}
\author{Astra P3}\date{27 October 2026}
\begin{document}\maketitle
\begin{abstract}
We construct exact counters that acquire hidden chain algebras from raw Boolean cubic coefficients. The strongest theorem is polynomial uniformly in unbounded module rank, residue degree and nilpotent depth on a certified family. A semilinear certificate manufactures repeated quadratic fibers, and eight exact moments reconstruct their joint contribution; each nontrivial moment is a cubic point count in at most eight field variables. The theorem uses the established Lauder--Wan backend, which is not implemented here. The executed reference terminal explicitly scans one field coordinate. Arbitrary higher nilpotent cubic coefficients and arbitrary-rank compatible quotient quadratics are allowed. A broader acquired contact family has $2^\kappa\operatorname{poly}(n,L)$ ring cost, with the binary kernel dimension $\kappa$ manufactured and charged; only the origin recurses. An exact reduction proves \#P-hardness when arbitrary kernels are allowed, even with fixed residue field, fixed depth and no error. A separate rank-three theorem covers every homogeneous cubic with unit mixed coefficient. Global low-rank errors, primitive products and split-cubic linear shifts are retained. Raw-input executions reach 100 bits, including an actual 72-input gate program and an 80-bit full terminal regression. All descriptions, arithmetic and reconstruction are measured against original input bits. Historical priority and practical speed advantage are not established; no generic circuit counter or polynomial SAT solver is claimed.
\end{abstract}

\section{The computational statement}
The input is $n$ and an explicit Boolean algebraic normal form
$p(x)=\bigoplus_{|S|\le3}a_S\prod_{i\in S}x_i$.
We charge its full description length $L$: index-list encoding uses
$O(n+(M+1)\log(n+2))$ bits for $M$ monomials; executable bit masks use
$O(n(M+1))$ bits. The exact signed and zero counts are
\[
 Z(p)=\sum_{x\in\F_2^n}(-1)^{p(x)},\qquad N_0(p)=(2^n+Z(p))/2.
\]
The number $n$ counts original independent bits, not gate states.
Let $F=\F_{2^d}$, $A=F[t]/(t^e)$, and
$\lambda(a)=\Tr_{F/\F_2}([t^{e-1}]a)$.
The first statement has $n=3ed$; the unbounded-rank extension has $n=bed$.
\subsection{The complete rank-three acquired family}
Let $H\in A[X,Y,Z]$ be any ordinary homogeneous cubic whose coefficient
$\mu=[XYZ]H$ is a unit. Let $Q$ be an arbitrary binary quadratic function
of $A^3/(t^hA)^3$, where $3h\le e+1$. The input may be any invertible
binary affine change of coordinates of
\begin{equation}\label{eq:quotient-family}
 p(x)=\lambda(H(x+a))+Q(x\bmod t^h).
\end{equation}
Neither $A,d,e,h,H,Q,a$ nor a useful coordinate map is supplied.

\begin{theorem}[Acquired weighted critical-ray counter]\label{thm:quotient}
For $d\ge3$, the family \eqref{eq:quotient-family} has a deterministic
polynomial-time exact counter from raw Boolean ANF. Every accepted
representation is verified by coefficient identities. There is no residue
smoothness assumption, factorization promise, or bound on the rank of $Q$.
One conservative bit-time bound is
\[
 O(Ln^2+n^8)+\widetilde O(d^{19}),
\]
with space $O(L+n^4)+\widetilde O(d^{12})$; the tilde terms occur only at
a depth-one finite-field terminal. The original number of independent
bits is $n=3ed$. All integer reconstruction uses $O(n)$-bit numbers.
\end{theorem}
The finite-field terminal uses a published explicit point-counting
algorithm, not an assumed oracle (Section~\ref{sec:fieldterminal}).
The accompanying generic implementation of that terminal is deliberately
limited to explicit enumeration within a stated bit budget. It never
labels that enumeration polynomial. For $e\equiv0,2\pmod3$ the executed
compiler needs no such terminal at any residue degree.

The quotient can contain $3dh$ independent bits and its quadratic polar
rank can reach the largest even integer below or equal to $3dh$.
For $e=2$ this is half of all original bits. As $e$ grows it approaches
one third. The count still uses no Walsh leaves. The coefficient-level
counter also works for $d=1,2$; completeness of this raw quotient fitter
is asserted for $d\ge3$. A separate smooth-residue compiler handles
those small degrees by exact derivative equations.

\subsection{The extension with unbounded module rank}
Section~\ref{sec:contact} constructs a counter for odd $b\ge5$ with no
fixed bound on $b$. The input is again an arbitrary binary affine disguise
of a homogeneous ring cubic plus an arbitrary quotient quadratic of height
$3h\le e+1$. The residue third form must admit a factor
$T_0=\ell\wedge\omega$, with $\omega$ nonsingular on $\ker\ell$.
The compiler acquires this factor by an exterior linear system and checks
the complete identity. In the acquired chart its gradient is
$Mw^{[2]}+Bw+c$; it verifies nonsingularity of $M$ and $B$ and computes
the binary kernel dimension $\kappa\le b-1$.

\begin{theorem}[Acquired counter at unbounded module rank]\label{thm:contact}
For $d\ge3$ and $e\ge2$, the preceding family has deterministic exact
counting time
\[
 \bigl(2^\kappa+\mathbf1_{e\equiv1\ (3)}2^d\bigr)\poly(n+L)
\]
and polynomial working space, with $n=bed$ original input bits.
All algebra, phase, center, chart and kernel data are acquired from raw
ANF. The compatible quadratic has no rank bound.
For each odd $b\ge5$ and each $d\ge3$ with $3\nmid d$, an explicit
residue family has $\kappa=0$, including after arbitrary higher nilpotent
cubic deformations. It is polynomial uniformly in $b,d,e$ for
$e\equiv0,2\pmod3$, and polynomial at every depth when $d=O(\log n)$.
\end{theorem}
This removes the fixed-module-rank obstruction of a growing Macaulay
certificate. The $2^d$ terminal is implemented as one-coordinate Gaussian
fibers, not as enumeration of all $2^{bd}$ field assignments.
The theorem is conditional on acquired, exactly checked structure, not
on a supplied representation. Section~\ref{sec:contact} also gives a
\#P-hardness reduction within the same geometric class when $\kappa$
is unrestricted. It does not assert that every large-kernel instance is hard.

\subsection{All-depth polynomial closure by eight acquired summaries}
A further semilinear certificate, proved in Section~\ref{sec:eightmoments},
manufactures the repeated-pair residue form \eqref{eq:pairnormal}.
It includes the one-ray family above and permits two acquired quadratic
pair types, an arbitrary nonzero phase scalar, and a radial cubic.
Neither its binary structure nor its two-dimensional blocks are supplied.
For $3\nmid d$, Theorem~\ref{thm:eightmoments} gives
\[
 O(Ln^2+n^{12})+\widetilde O(d^{31})
\]
bit time at every depth, uniformly in unbounded $b,d,e$.
Seven degree-three point counts in at most eight variables and one
binary Gaussian manufacture eight linear summaries. Exact reconstruction
then raises their acquired states to the numbers of pair types.
The generic point-count engine is a published dependency, not an executed
component here; the reference tests explicitly record their field scans.
The degree restriction is substantive: Section~\ref{sec:eightmoments}
also embeds arbitrary Boolean cubic counting into this same repeated-pair
residue at $d=3,e=2$, using only higher nilpotent coefficients and no
quadratic noise.

\subsection{The complementary global-error family}
The base family consists of invertible binary affine disguises of
\begin{equation}\label{eq:family}
 \lambda(\mu XYZ+A_1XY+A_2XZ+A_3YZ+B_1X+B_2Y+B_3Z+B_0),
 \qquad \mu\in A^\times.
\end{equation}
Other coefficients are arbitrary ring elements, including zero divisors.
An additional binary quadratic error is allowed. The compiler receives none
of this representation, its coefficients, or its error and rank.

After acquisition, view the space as $F^{3e}$. Let $G$ be the block matrix
of the ordinary field trace pairing. If $D$ is the error's alternating polar
matrix, put $E=G^{-1}D$, $U=\im E$ and $r=\dim_{\F_2}U$.
Section~\ref{sec:sylvester} manufactures a seven-dimensional scalar space
$S\subset F$ for every $d\ge7$. The sufficient recovery condition is
\begin{equation}\label{eq:direct}
 \dim_{\F_2}(SU)=7r.
\end{equation}
It is checked on a recovered support, not supplied as a decomposition.

\begin{theorem}[Complete acquired local-algebra counter]\label{thm:main}
For $d\ge7$, every input above whose actual support satisfies
\eqref{eq:direct} is counted deterministically in $2^r\poly(L+n)$ bit time
and polynomial space. There is no bound on $e$ and no requirement $r<d/2$.
Returned answers have exact coefficient identities; failed structural tests
cause rejection. For independent uniform support columns the failure
probability is less than $2^{7r-n}$; for a uniform rank-$r$ alternating
polar matrix it is at most $2^{7r-n+1}$.
\end{theorem}
The unperturbed result holds for every $d\ge3$. The previous Gabidulin
decoder gives a second, worst-case route when $2r\le d-3$. Adversarial
supports can favor that older route.

For fixed $d\ge7$ and $e\to\infty$, residue-field descent alone creates
$3e$ field variables. A fixed-variable point counter ceases to give a
polynomial bound. The nilpotent evaluator removes that obstruction,
while global support recovery removes the radius tied to $d$. The new
bound is polynomial for $r=O(\log n)$ and subexponential for $r=o(n)$
on the stated family. It is not a generic circuit counter.

\section{Acquiring the full local algebra}
The alternating third difference $T(u,v,w)=\Delta_u\Delta_v\Delta_w p$
and its radical are explicit. Compute the balanced input-slot centroid
\[
 \Cent(T)=\{C:T(Cu,v,w)=T(u,Cv,w)=T(u,v,Cw)\}.
\]
There are $n^2$ matrix unknowns and $O(n^3)$ binary equations.
For radical-free $T$ the centroid is commutative: moving two operators
between three slots puts their commutator in the radical. Multiplication
and the identity are verified in the computed algebra.

For \eqref{eq:family} this centroid is exactly $A$. Balancing against vectors
with two equal module positions forces off-diagonal blocks to vanish.
Nondegeneracy of $\lambda(ab)$ then gives $L_1(x)y=xL_2(y)$ for diagonal
blocks. Setting $x=1$ or $y=1$ makes every block multiplication by a common
element. The unit $\mu$ preserves nondegeneracy. Binary affine disguise
conjugates the centroid, and a quadratic error leaves $T$ unchanged.

\begin{lemma}[Constructive chain presentation]
A finite commutative binary matrix algebra can be recognized as
$\F_{2^d}[t]/t^e$, and this presentation acquired, in deterministic
polynomial time. Freeness of its input-space module can also be recognized,
and a free module basis acquired in polynomial time.
\end{lemma}
\begin{proof}
Squaring $\sigma:c\mapsto c^2$ is binary linear. The test
$\dim\ker(\sigma-I)=1$ recognizes a local algebra, since its fixed elements
are idempotents. Iterated squaring stabilizes its image. The stable image
is the canonical coefficient field $F$; the kernel of a sufficiently high
iterate is the nilradical $J$. Iteration kills nilpotents while inducing
an automorphism on the finite residue field.

Compute $J^2$ and test $\dim_{\F_2}(J/J^2)=d$. Any
$\pi\in J\setminus J^2$ generates $J$ by Nakayama's lemma.
If $\dim C=de$, explicitly verify $\pi^e=0$ and that
$\{\alpha^i\pi^j:0\le i<d,\ 0\le j<e\}$ is a basis. This certifies
the entire chain presentation and multiplication table.

To acquire $\alpha$ deterministically, factor $d$ into coprime prime
powers $m=\ell^a$. Frobenius kernels find an element of degree $m$
outside the unique maximal proper subfield of $\F_{2^m}$. Their sum
has degree $d$: a Frobenius displacement fixing a sum of two coprime-degree
elements lies in their intersection $\F_2$, and its iteration over the
odd-degree factor forces that displacement to vanish. Induct on the
prime-power factors. The matrix minimal polynomial presents $F$; for
$d=1$ take $\alpha=1$.

The nilpotent operator $\pi$ has blocks of length at most $e$.
Its module is free exactly when $e\dim\ker\pi=n$.
Acquire an $F$-basis of $V/\pi V$, lift it, and take its
$\alpha^i\pi^j$ multiples. Nakayama and the dimension test certify an
invertible binary coordinate matrix.
\end{proof}

The pairing $\lambda(ab)$ is nondegenerate: an element of valuation $j$
pairs nontrivially with a suitable multiple of $t^{e-1-j}$.
Thus $T$ has a unique $A$-multilinear lift defined by
\[
 \lambda(c\widetilde T(u,v,w))=T(cu,v,w),\qquad c\in A.
\]
Balancing proves multilinearity and alternation. On a free rank-three
module it is a volume form with unit coefficient $\mu$.
Therefore $p-\lambda(\mu XYZ)$ is exactly quadratic.
The compiler solves every quadratic and affine coefficient; it need not
reevaluate the cubic truth table to reprove the balancing identity.
An independent finite-test option checks every transformed third coefficient.

\section{Acquiring a homogeneous phase modulo quotient noise}\label{sec:quotientfit}
This section proves the acquisition part of Theorem~\ref{thm:quotient}.
The full third tensor of \eqref{eq:quotient-family} is the scalar volume
form with unit coefficient $\mu$. Repeated-variable ordinary cubic
monomials, including $X^3$, have binary degree at most two in
characteristic two. Thus all third coefficients determine exactly the
term $\lambda(\mu XYZ)$, while the remaining Boolean quadratic is explicit.
The preceding centroid construction therefore acquires $A$ and its
rank-three free module without a supplied ring presentation.

\subsection{Why the promise survives acquired coordinates}
Two presentations of this full centroid differ by a ring isomorphism
$\sigma$ and an invertible module matrix. Including translations, the
change has the form $x=M\sigma(y)+b$, with $M\in\GL_3(A)$.
Nondegeneracy of the Frobenius pairing gives a unique unit $c$ such that
the transported functional is $z\mapsto\lambda'(cz)$. Consequently
the transformed ordinary homogeneous cubic is
$c\sigma^{-1}(H\circ M)$, with a translated center.
Its mixed coefficient is
$c\sigma^{-1}(\mu\det M)$, again a unit. A repeated-variable cubic
cannot create an $XYZ$ term under linear substitution in characteristic two.
The ideal $J^h$ is preserved by both $M$ and $\sigma$.
The transported $Q$ is still a binary quadratic on $A^3/(J^h)^3$.
This argument includes nonlinear changes of uniformizer and coefficient-field
Frobenius; no preferred hidden basis is assumed.

\begin{lemma}[Ordinary coefficients invisible above a quotient]\label{lem:invisible}
Let $d\ge3$. Suppose the Boolean function $\lambda(P(x))$, where $P$ is
ordinary of total degree at most three, depends only on $x\bmod J^h$.
After replacing ordinary square terms by their exact linear adjoints,
every homogeneous cubic coefficient and every cross-quadratic
coefficient of $P$ lies in $J^{e-h}$.
\end{lemma}
\begin{proof}
For all $x$ and $v\in(J^h)^3$, the difference
$\lambda(P(x+v)+P(x))$ is zero. Scale both $x$ and $v$ by every
$\rho\in F$. Writing the homogeneous pieces of $P$ separately gives
\[
 \Tr\bigl(a\rho^3+b\rho^2+c_1\rho\bigr)=0
 \quad\text{for every }\rho\in F,
\]
where $a=[t^{e-1}](P_3(x+v)+P_3(x))$.
For $d\ge3$, the doubling orbit of exponent $3$ modulo $2^d-1$ has
length $d$ and is disjoint from the orbit of $1$ (which contains $2$).
Hence the trace polynomial identity forces $a=0$.
This is an ordinary polynomial identity over $F$ in the coefficient
variables of $x$ and $v$; its degree is at most three and $|F|\ge8$,
so its coefficients vanish.

Each degree-three monomial has an odd exponent. To isolate its coefficient
at $t^j$, for $j<e-h$, use once the variable $v_i$ at power
$t^{e-1-j}$ in such an odd-exponent slot, and use residue variables
for the other occurrences. The binomial coefficient is odd and no other
monomial contributes the same coefficient-variable monomial. Thus that
coefficient is zero. This proves $P_3\in J^{e-h}A[X,Y,Z]$.
It already defines a function on the low quotient, so subtract it.
Binary polarization between different module coordinates and
nondegeneracy of the field trace pairing then force every cross-pair
coefficient into $J^{e-h}$ as well. Square terms have been folded into
linear terms and do not enter this last polarization. Once the cubic and
cross terms are removed, the same Frobenius pairing forces every remaining
linear coefficient into $J^{e-h}$ as well.
\end{proof}

\subsection{A linear fit that does not guess the error}
Set $h=\lfloor(e+1)/3\rfloor$; this includes every smaller allowed quotient.
Keep the acquired term $\mu XYZ$ fixed. For the other nine ordinary
homogeneous cubic monomials, three cross pairs, three linears and a
constant, form their binary trace-ANF columns, one for each ring
coefficient bit. There are $16ed$ columns. All have binary degree at
most two, so every acquired third coefficient remains fixed.
Constrain every quadratic or linear ANF coefficient involving a bit
outside the low quotient. The remaining rows may be arbitrary noise.
The protected coefficient matrix has exact rank $15d(e-h)$: its kernel
consists of arbitrary constant coefficients and the $15$ nonconstant
coefficient families in $J^{e-h}$, hence has dimension $de+15dh$.
Solve this binary linear system and obtain any ordinary polynomial $P$.
Failure rejects the input.

Let $H=P_3$ and read the affine center directly:
\[
 a=(P_{YZ},P_{XZ},P_{XY})/\mu.
\]
Subtract $\lambda(H(x+a))$ from the original acquired phase. Its third
tensor is zero, so all remaining coefficients are recovered exactly
from evaluations of weight at most two. Verify that every resulting
monomial is supported on the low quotient. Finally translate $x$ by
$a$ and carry the translated quotient quadratic $Q'$ explicitly.
The output is the exact identity
\[
 p(x+a)=\lambda(H(x))+Q'(x\bmod J^h).
\]

To prove completeness, compare any fitted $P$ with the ordinary trace
polynomial obtained from the promised translated homogeneous phase.
Their trace difference depends only on the quotient. By
Lemma~\ref{lem:invisible}, their homogeneous cubics and cross-pair
coefficients differ by elements of $J^{e-h}$. Their centers therefore
differ by $\delta\in(J^{e-h})^3$. Substituting such a shift in a cubic
changes its trace only through the low quotient: coefficients of every
new nonconstant term contain a factor from $J^{e-h}$, while terms
with two such factors vanish because $e\ge2h$.
Thus every linear-system solution produces an admissible normalized
identity. This removes a search among fit solutions.
The final coefficient check also certifies accepted inputs outside the
promised generation history.

\section{Exact counting by manufactured critical rays}\label{sec:criticalrays}
Now $H$ and $Q$ are acquired explicitly. No factorization of $H$ is
needed. Write $q=2^d$, $\psi_e(z)=(-1)^{\lambda(z)}$, and let
$H_0$ denote reduction modulo $t$.
For $e\ge2$, averaging over the highest digit of each variable cancels
every nonzero residue vector where $\nabla H_0\ne0$. The quotient
quadratic is independent of those highest digits. The surviving
nonzero residue classes lie on projective singular rays of $H_0=0$.
Euler's degree-three identity in characteristic two gives
$H_0(x)=x\cdot\nabla H_0(x)$, so a gradient zero lies on the curve.

\subsection{Finding every ray by binary linear algebra}
Use disjoint projective charts: $Z\ne0$ first; then $Z=0,Y\ne0$;
then $Z=Y=0,X\ne0$. In the first chart put $Z=1$ and
$f(u,v)=H(u,v,1)$. Its gradients have the form
\[
 f_u=a u^2+b v^2+\mu v+c_1,\qquad
 f_v=c u^2+d_1 v^2+\mu u+c_2.
\]
Over the residue field this is an affine binary linear system with
$2d$ unknown bits. Its Jacobian is
$\left(\begin{smallmatrix}0&\mu_0\\\mu_0&0\end{smallmatrix}\right)$,
with nonzero determinant. Two quadratics with this Jacobian have at
most four geometric solutions; hence the binary solution space has
kernel dimension at most two. Enumerate these at most four candidates,
apply the chart restrictions, and retain $H_0=0$.

There are at most three projective singular rays. A repeated plane factor
would make $H_0=L^2M$ and annihilate its $XYZ$ coefficient. Thus the
cubic is reduced. A line through two singular points divides a cubic,
since the two intersection multiplicities already total at least four.
Three distinct singular points therefore force a triangle of lines;
there cannot be a fourth. This also covers nonsplit residue curves.

Each transverse critical point has a unique lift $a\in A^2$, because
the displayed Jacobian is a unit. At layer $t^k$ solve the two residual
gradient coefficients using its inverse. All operations are in the
acquired ring. The lifted vector $v=(a_1,a_2,1)$, the critical value
$c=f(a)$, and the residue quadratic coefficients below are manufactured.
Their truncations are the unique lifts at every smaller depth, so they
can be reused throughout the origin recursion.

\subsection{The exact transverse calculation in characteristic two}
Translate to the lifted critical point. The chart phase becomes
\[
 f(a+w)=c+A w_1^2+\mu w_1w_2+B w_2^2+C_3(w_1,w_2),
 \qquad w\in(tA)^2,
\]
where $C_3$ is homogeneous cubic. Let $s=\lfloor e/2\rfloor$ and
$k=\lceil e/2\rceil$. Partition the transverse domain into cosets of
$(t^kA)^2$. Addition of $\delta\in(t^kA)^2$ changes the phase exactly
by its formal gradient paired with $\delta$, because $\delta^2=0$ in
$A$. The weight $Q$ is unchanged, since $h\le s\le k$.
Character orthogonality therefore forces
$\nabla f(a+w)\in(t^sA)^2$.

The gradient is an invertible linear term $\mu(w_2,w_1)$ plus a
quadratic-in-$w$ term. On $(tA)^2$ it preserves the minimum valuation:
for a nonzero vector of valuation $j\ge1$, the linear term has
valuation $j$ and the quadratic term at least $2j$.
Thus the forced condition is exactly $w\in(t^sA)^2$.
For even $e=2s$, all nonconstant terms then vanish and the transverse
sum is $q^e$. For odd $e=2s+1$, the highest free digits contribute
$q^{2s}$ and the two middle digits contribute the field Gaussian
\[
 \sum_{u,v\in F}(-1)^{\Tr(A_0u^2+\mu_0uv+B_0v^2)}
   =q\kappa,\qquad
 \kappa=(-1)^{\Tr(A_0B_0/\mu_0^2)}.
\]
The cubic term vanishes because $3s\ge e$ for odd $e\ge3$.
For a fixed radial unit $r$, every coefficient is multiplied by $r^3$,
which cancels from the displayed Arf ratio. The quotient weight is
constant on this surviving transverse cell and equals $(-1)^{Q(rv)}$.
The exact contribution of one projective ray is consequently
\begin{equation}\label{eq:weightedray}
 q^e\kappa^{e\bmod2}
 \sum_{r\in A^\times}(-1)^{\lambda(c r^3)+Q(rv)}.
\end{equation}
No characteristic-zero Morse lemma is used.

\subsection{A quadratic radial phase and a single recursive state}
The map $r\mapsto r^2$ is binary linear. Hence
$\lambda(c r^3)+Q(rv)$ is a binary quadratic on $ed$ bits, even when
$Q$ has large rank. The unit sum in \eqref{eq:weightedray} is exactly
\[
 \mathcal G_A\bigl(\lambda(c r^3)+Q(rv)\bigr)
 -\mathcal G_{tA}\bigl(\lambda(c r^3)+Q(rv)\bigr),
\]
where each $\mathcal G$ is computed by ordinary binary symplectic
elimination, including its radical and affine term. The two Gaussian
variable counts are $ed$ and $(e-1)d$. No character expansion of $Q$
is formed. For $Q=0$, the unit sum also has the direct formula
\[
 U_e(c)=\begin{cases}
 0,&v_t(c)\le e-2,\\
 q^{e-1}\left(\sum_{r\in F}(-1)^{\Tr(br^3)}-1\right),
       &c=t^{e-1}b.
 \end{cases}
\]
This formula includes $c=0$. The field cube phase is binary quadratic;
no claim that cubing permutes $F^\times$ is needed.

The only remaining residue class is the origin. Put $x=ty$.
For $e\ge4$, homogeneity gives the exact origin contribution
\[
 q^6 Z_{e-3}\bigl(H\bmod t^{e-3},\ Q(ty)\bigr).
\]
There are two unused ring digits in each of three variables. The
quotient height decreases by one, and
$3(h-1)\le(e-3)+1$, so the same hypothesis persists.
For $e=2$ or $3$, the homogeneous cubic vanishes on the origin class
and the contribution is
$(-1)^{Q(0)}q^{3(e-1)}$.
At most three nonrecursive ray contributions are added at each level.
Only one origin state recurses, for at most $\lceil e/3\rceil$ levels.
This is why the number of computations is polynomial rather than an
exponentially branching lift tree.

\subsection{The finite-field terminal}\label{sec:fieldterminal}
When the origin recursion reaches depth one, the allowed quotient
height is zero. The remaining phase is
$\Tr(H_0(X,Y,Z))$ plus a constant. Manufacture the ordinary equation
\[
 V:\quad W^2+W+H_0(X,Y,Z)=0
\]
over $F$. Each triple has two lifts $W$ when its trace is zero and
none otherwise, so
\[
 Z_1=\#V(F)-q^3.
\]
This is a four-variable, degree-three, characteristic-two point count.
Lauder--Wan, Proposition~36 and the affine torus decomposition in
Section~6.4, give deterministic bit time $\widetilde O(d^{19})$ and
space $\widetilde O(d^{12})$ for this first point count \cite{lauderwan}.
The constants depend only on the fixed degree, dimension and characteristic;
no smoothness is required. Computing an entire zeta function is unnecessary.

The generic point-count engine is not implemented in this source package.
The executed code explicitly limits its field terminal to $3d\le24$
bits by default, and records every enumerated point. For
$e\equiv0,2\pmod3$, it never reaches a field terminal and has no
field-size restriction. The separate file-leaf compiler manufactures
and verifies the full trace-polynomial identity using $20d$ binary
coefficient columns. It rejects arbitrary binary quadratics that do
not have an ordinary degree-three field lift. Its Artin--Schreier
reduction is executed independently; it is not advertised as an
execution of the general Lauder--Wan algorithm.

An additional implemented specialization handles the field family
$X^2Y+XY^2+Z^3+aXYZ+\rho X+\sigma Y+\tau Z+c$.
Eliminating $X$ and substituting $Y=w^2$ leaves a Laurent trace phase
with exponents among $\{3,1,-1,-3,-5\}$, plus an exact Gaussian
boundary contribution. The finite Dwork trace backend has at most
$32(d+2)-1$ states and conservative $O(d^8)$ bit time and $O(d^4)$
space. This is an executed specialization of established Dwork machinery.
No acquired equivalence from every cubic to that special form is claimed.

\subsection{Description size and arithmetic costs}
The acquired binary coordinate map uses $n^2$ bits. Ten ring coefficients
use $10ed$ bits. A quotient quadratic uses at most
$1+3dh+\binom{3dh}{2}$ binary coefficients, or $O(n^2\log(n+2))$
bits in an index-list encoding. The complete representation is therefore
polynomial in the original independent-input count. It is manufactured
from all raw coefficients, not passed to the counter as advice.

For an explicit conservative cost, binary centroid elimination has
$n^2$ unknowns and $O(n^3)$ equations, and costs at most $O(n^7)$
bit operations with classical elimination. Ring structure, Frobenius
subspaces and free module bases use polynomial-size binary matrices;
the algorithms described above fit within $O(n^8)$ bit time.
The quotient fit has $O(n)$ columns and $O(n^2)$ rows. Each candidate
trace-ANF is obtained from ring monomial evaluations and a Gram matrix;
all ring elements have at most $n$ bits. Its construction, solve,
normalization and exact coefficient checks fit within the same bound.

There are $O(e)$ stages and at most six binary Gaussian sums per stage.
Each has at most $ed\le n$ variables. Even constructing its coefficients
by all weight-at-most-two evaluations, each evaluation scanning an
$O(n^2)$-term quadratic with $n$-bit masks, gives a bound below $O(n^8)$
for the entire recurrence. Schoolbook multiplication of ring coefficients
and of count integers is included. Streaming the binary equations and
retaining polynomial-size matrices gives $O(L+n^4)$ space before a
finite-field terminal. Combining these bounds with the displayed
Lauder--Wan term proves Theorem~\ref{thm:quotient}.

All count contributions and recursion multipliers have $O(n)$ bits.
The original signed sum lies between $-2^n$ and $2^n$; the final zero
count is reconstructed as $(2^n+Z)/2$. No unit-cost arithmetic over
exponentially large integers is assumed. For an accepted gate program,
add the explicit gate-to-ANF cost from the preceding section; internal
wire values are never counted as independent inputs.

\section{Smooth shortcuts, sharp failures, and an adaptive repair}
For a smooth residue cubic there are no nonzero critical rays. The
weighted recurrence therefore becomes exact digit forcing toward a
single translated origin. For a homogeneous centered phase and the
allowed quotient height,
\[
 Z_e(H,Q)=(-1)^{Q(0)}
 \begin{cases}
 q^{2e},&e\equiv0\pmod3,\\
 q^{2e-1},&e\equiv2\pmod3,\\
 q^{2(e-1)}Z_1(H_0),&e\equiv1\pmod3.
 \end{cases}
\]
The earlier smooth-only compiler manufactures the residue center directly
from highest-digit derivatives. After fitting each binary-linear map as
$M_{ij}X_j^2+B_{ij}X_j$, it reads the center
\[a=(B_{23},B_{13},B_{12})/\mu_0,\]
checks every gradient identity, and
certifies absence of nonzero homogeneous gradient zeros with an
$18\times15$ Macaulay matrix in degree four. Full rank means that the
three quadrics generate every quartic monomial. Conversely a zero-free
homogeneous gradient is a regular sequence of quadrics, whose quotient
has Hilbert series $(1+t)^3$, so this certificate is complete.
For $d=2$ the square and linear maps are still distinguishable; for $d=1$
the implementation checks the exact eight residue vectors instead.

The smooth-only implementation also checks that the proposed discarded
layers are absent. This is essential: over $\F_4[t]/t^4$ with
$H=X^2Z+Y^3+Z^3+XYZ$, adding one linear bit from the $t^2$ coefficient
leaves the first center and smoothness test unchanged but changes the
signed count from $114688$ to $0$. The unused-layer test rejects the
invalid compression. The independent audit executed both full truth tables.

The sign-only invisibility identity has a sharp height limit. Over
$\F_2[t]/t^e$ use the same smooth $H$ and set
$h_0=\lfloor(e+1)/3\rfloor$. A rank-two quadratic using the next digit
$t^{h_0}$ changes $Z$ by a factor $1/2$, $-1$, or $2$ according as
$e\equiv0,1,2\pmod3$: use $X_{h_0}Y_{h_0}$ in the first and third
cases and $X_{h_0}Z_{h_0}$ in the second. The exact terminal signed
sums at depths $3,1,2$ are respectively $64,2,8$, and become
$32,-2,16$. Valid origin reductions propagate these examples to all
depths. This refutes the tempting simplification, not polynomial counting
outside the sufficient height range.

Indeed the critical-ray calculation only needs $2h\le e$ at the current
stage. Continue it while this holds, then keep the single remaining
state explicitly. If
$\Delta=h-\lfloor(e+1)/3\rfloor>0$, the residual number of binary
variables is at most
\[
 R\le\min\{n,\ 3d\max(1,6\Delta-1)\}.
\]
At level $j$ the depth is $e-3j$ and the nonzero height is $h-j$;
the condition is $j\le e-2h$. This gives the displayed bound on the
first unhandled state. The executable accepts a separate residual-bit
budget and records its $2^R$ enumeration explicitly. The additional
cost is $2^R\poly(L+n)$, not exponential in the full rank of $Q$.
For fixed $d$ and fixed excess height, it is polynomial in growing $e$.

The raw fitter can also acquire the minimum admissible height. The
coefficient-ambiguity lemma holds for every $0\le h\le e$. Its center
argument does not actually need $\delta^2=0$: every changed coefficient
contains $\delta\in J^{e-h}$ and hence remains supported on the low
quotient, even when higher powers of $\delta$ survive.
Existence of the fitted representation is monotone in $h$, and height
$e$ always permits $H=\mu XYZ$ plus the full explicit quadratic.
Binary search using the exact fit therefore manufactures the smallest
height. The semilinear invariance argument makes this minimum independent
of the acquired chain presentation. No favorable height is
supplied. Combining it with the residual bound gives a certified
adaptive algorithm on every raw rank-three chain-volume cubic with
$d\ge3$, with graceful but explicitly exponential cost in an unfavorable
residual. It does not turn arbitrary dense quadratic couplings into a
polynomial family.

\section{Obstructions that the construction does not erase}
\subsection{Arbitrary linear constraints recover general cubic counting}\label{sec:parityboundary}
The unrestricted character sum of a ring or field phase can be easy while
its restriction to an arbitrary binary subspace is hard. This matters for
conditional queries, inversion, synthesis and autoregressive learning:
none follows from a total count without checking the new constraints.

Here is an explicit reduction. For any Boolean cubic $p$ on $m$ bits,
choose $F=\F_{2^{m^2}}$ with polynomial basis $1,\alpha,\ldots$.
Put $a_i=\alpha^i$ and $b_j=\alpha^{mj}$ for $0\le i,j<m$.
Their $m^2$ products are a basis. For each $k$, solve the nondegenerate
trace-pairing system for $c_k$ so that
$\Tr(a_i b_j c_k)$ is an assigned coefficient array. Assign each Boolean
monomial of $p$ to one ordered triple, permitting repeated indices for
linear and quadratic terms. Then the injective binary map
\[
 u\longmapsto\left(\sum_i a_i u_i,\sum_j b_j u_j,\sum_k c_k u_k\right)
\]
restricts $\Tr(XYZ)$ to $p$ (a separate constant sign is harmless).
A parity-check description of its image is manufactured by elimination
and uses $O(m^4)$ bits. The fixed-characteristic field construction is
polynomial; standard deterministic constructions are available \cite{shoup}.

Cubic signed counting is already \#P-hard: encode every Boolean circuit
gate by its unique-extension equation $f_j=0$ of degree at most two,
including an accepting-output equation, and introduce independent bits
$z_j$. Orthogonality gives
\[
 \sum_{x,\text{gate states},z}(-1)^{\sum_j z_j f_j}
   =2^{\#\{j\}}\#\{\text{accepting original inputs}\}.
\]
Thus a general polynomial counter for arbitrary parity-constrained
$\Tr(XYZ)$ would imply $\mathrm{FP}=\#\mathrm P$.
This is a reduction, not an assumed separation.

A separate Fourier calculation shows that a rank-$r$ quadratic sign has
exactly $2^r$ nonzero Walsh coefficients. Any universal reconstruction
using fixed linear-character shifts and weights needs that many terms.
This is a restriction on that representation, not a lower bound on
counting time. The weighted critical-ray method explicitly escapes it
for compatible quotient weights by using a different reconstruction.

\subsection{A nonchain centroid can retain an arbitrary tensor problem}
For any symmetric trilinear form $S$ on $F^k$, construct the commutative
Frobenius algebra
\[
 A_S=F\oplus U\oplus U^*\oplus F\omega,\qquad U=F^k.
\]
Define $uv=S(u,v,\cdot)\in U^*$, $u\phi=\phi(u)\omega$, and all other
positive-degree products by the grading. Symmetry of $S$ gives
associativity, the radical has fourth power zero, and the functional
extracting $\omega$ has a nondegenerate multiplication pairing.
The centroid of the scalar volume tensor on $A_S^3$ is exactly $A_S$.
Nevertheless its cubic signed sum satisfies the explicit identity
\[
 Z_{A_S}=q^{3k+3} Z_S+3(q-1)q^{4k+3},
\]
where $Z_S=\sum_{u,v,w\in F^k}(-1)^{\Tr S(u,v,w)}$.
Thus acquisition of an arbitrary local algebra can leave the original
trilinear counting problem inside it. The polynomial chain-ring theorem
cannot be extended merely by replacing the word ``chain'' with ``local''.
The algebra construction and the displayed identity have complete
proofs and executed multiplication/count checks in the source files.

\subsection{Other preserved failures}
A residue cubic may split into three independent planes without its full
ring lift splitting: the pure-cube coefficients in the Hensel residual
are genuine obstruction checks. The global nilpotent annihilator may
annihilate a defect in its kernel even though the residual is not an
ordinary ring cubic; the final coefficient fit is essential.
Small residue fields can identify ordinary cubic trace coefficients,
which is why the quotient-fit completeness argument retains $d\ge3$.
The square adjoint in \eqref{eq:square} must include nilpotent indices;
a field-only square-root formula is false.

Finally, the new quotient family permits arbitrary-rank compatible
quadratics, but not arbitrary ordinary linear perturbations. For example,
$\lambda(XYZ+X)$ at depth two fails the translated-homogeneous low-quotient
fit. It remains supported by the complementary split-cubic counter.
This failure is retained as an actual rejection example, not bypassed
by dropping unmatched coefficients.


\section{A fixed-rank extension beyond ternary modules}\label{sec:rankfive}
The unique-origin mechanism extends beyond three ring coordinates.
Here write $n=bed$, where $b$ is the acquired free-module rank.
Suppose the full balanced centroid is a chain algebra, its module is free,
and the input is an affine binary disguise of
\[
 \lambda(H(X_1,\ldots,X_b))+Q(X\bmod t^h),\qquad 3h\le e+1,
\]
where $H$ is ordinary homogeneous cubic and its residue gradient has no
nonzero zero over the algebraic closure. The acquired third tensor is
required to have zero radical; the algorithm checks this condition.
Smoothness alone would be insufficient for acquisition: a sum of pure
cubes can be geometrically smooth and still have zero binary third tensor.

Let the acquired squarefree ring cubic be
$C(X)=\sum_{i<j<k}C_{ijk}X_iX_jX_k$ and let the explicit remainder be
binary quadratic. Extend $C_{ijk}$ symmetrically, with repeated indices
set to zero. Highest-digit derivatives manufacture field quadrics
\[
 g_i(X)=\sum_{j<k}C_{ijk,0}X_jX_k+
        \sum_jM_{ij}X_j^2+\sum_jB_{ij}X_j+c_i.
\]
For $d\ge2$, binary coefficient fitting distinguishes the maps $X$ and
$X^2$. The translated center is the unique solution of
\[
 B_{ij}=\sum_k C_{ijk,0}a_k.
\]
The coefficient matrix is the contraction map of the residue third tensor;
its rank is checked to be $b$. Verify the remaining constant equations
$g_i(a)=0$.

For the homogeneous parts of the $g_i$, construct the Macaulay matrix
that multiplies each quadric by every degree-$(b-1)$ monomial. Its row
and column counts are
\[
 R_b=b\binom{2b-2}{b-1},\qquad N_b=\binom{2b}{b-1}.
\]
Full column rank certifies that every degree-$(b+1)$ monomial is in the
gradient ideal, hence there is no nonzero common zero. Conversely,
a zero-free homogeneous system of $b$ quadrics in $b$ variables is a
regular sequence, and its quotient has Hilbert series $(1+t)^b$;
there is no surviving degree $b+1$. Thus the rank certificate is complete,
not merely a sufficient random probe. At $b=5$ the matrix is $350\times210$.

Highest-digit orthogonality now forces the unique residue $a$ in all
$b$ variables. Restrict $X=a+tY$, manufacture its remaining quadratic,
and verify every coefficient of the two layers proposed for removal.
For $e\ge4$, their exact factor is $q^{2b}$ and the new depth is $e-3$.
At depths two or three the remaining phase is binary quadratic and is
eliminated exactly. For each fixed $b$, all construction and reconstruction
costs are polynomial in $n+L$. The Macaulay dimensions above must be
charged when $b$ varies; there is no polynomial claim for unbounded $b$.
An optional depth-one terminal reduces to a degree-three
Artin--Schreier equation in $b+1$ field variables. An exact binary linear
fit with $d\binom{b+3}{3}$ ordinary coefficient columns manufactures the
terminal polynomial and verifies its complete trace-ANF identity. The published terminal
bound is $\widetilde O(d^{3b+10})$ time and
$\widetilde O(d^{2b+6})$ space. The present executable requires an exact
callback there and never silently enumerates that field.

\subsection{An explicit rank-five family and its full centroid}
Use coordinates $(X_1,Y_1,X_2,Y_2,Z)$ and residue polynomial
\[
 H_0=Z^3+Z(X_1Y_1+X_2Y_2)+
 \sum_{i=1}^2\left(X_i^2Y_i+c_i^{-1}X_iY_i^2\right).
\]
Choose $c_1,c_2$ so that $1$ is not in their binary span; over $\F_8$
one may use $c_1=\alpha,c_2=\alpha^2$ in a polynomial basis.
A critical point with $Z=0$ is zero. With $Z=1$, the first four gradient
equations force $X_i\in\{0,1\}$ and $Y_i\in\{0,c_i\}$.
The last becomes $1+\varepsilon_1c_1+\varepsilon_2c_2=0$, impossible.
This proves geometric smoothness without scanning residue points.
Arbitrary higher nilpotent coefficients may be added to every ordinary
cubic monomial; they do not change this argument.

The residue third tensor is $Z$ wedged with a nondegenerate alternating
form on four variables. Its balanced centroid consists only of scalars:
balancing three vectors in the four-dimensional subspace removes the
component toward $Z$; balancing one $Z$ forces the remaining block to
be a common scalar; balancing two $Z$'s removes the reverse component.
Its radical is zero.

This family is nonempty at every fixed odd rank $b=2k+1\ge5$ and
every $d\ge2$. Use the same formula with $k$ pairs and choose every
$c_i=\alpha\in F\setminus\F_2$. The last gradient equation becomes
$1+\alpha\sum_i\varepsilon_i=0$, which is impossible because the
sum lies in $\F_2$. The third tensor is $Z\wedge\omega_{2k}$, and
the same scalar-centroid argument applies. The executions below are
for rank five; no higher-rank execution is claimed.

For the full scalar tensor over $A$, the algebra $A$ is already contained
in the balanced binary centroid. The latter is commutative because the
tensor is radical-free. Every centroid operator therefore commutes with
$A$ and is $A$-linear. Modulo $t$ it is a scalar by the preceding argument.
Subtract a scalar lift and divide by $t$; the balance equations descend
one level. Induction forces the operator to be multiplication by a single
ring element. Thus the full binary centroid is exactly $A$, even with
arbitrary higher mixed coefficients. The useful ring is genuinely acquired
from raw disguised coefficients.

The executed raw rank-five examples have $(n,d,e)=(30,3,2)$ and
$(45,3,3)$, each with a rank-$14$ quotient quadratic. Their signed counts
are respectively $32768$ and $-1073741824$. The counter acquires centroid
dimensions $6$ and $9$, obtains module rank five, and verifies the
$350\times210$ certificate. Each case also passed $1024$ independent
original-input representation identities. These are not instances
of the rank-three compiler under a relabeling of its original inputs.

Another raw case uses $d=2,e=2$ and $c_1=c_2=\alpha\in\F_4\setminus\F_2$;
the same smoothness argument applies. All $1048576$ assignments of this
$20$-bit instance, with rank-$10$ quotient noise, give $Z=-1024$ and
$N_0=523776$. A $75$-bit case with $d=3,e=5$ and rank-$30$ noise
verifies the actual recursive step $5\to2$ and factor $2^{30}$; its
signed count is $-35184372088832$. A separate $120$-bit coefficient-state
check verifies $8\to5\to2$ and is not described as raw acquisition.
Highest-digit linear, discarded-layer linear, and projectively singular
perturbations were each rejected by the corresponding exact test.


\section{An acquired counter without a fixed module-rank bound}\label{sec:contact}
The fixed-rank smoothness certificate is unnecessary on a larger acquired
family. Put $b=2r+1\ge5$, $n=bed$, and allow $b$ to grow.
Suppose the residue third form is $T_0=\ell\wedge\omega$, where
$\omega$ is nonsingular on $W=\ker\ell$. This identity is manufactured
and checked below; neither factor is supplied. In the chart $z=\ell=1$,
the transverse gradient is $g(w)=Mw^{[2]}+Bw+c$, with $B$ alternating
and nonsingular. Require $M$ nonsingular when $e\ge2$.
Define $\kappa$ as the computed binary nullity of $w\mapsto Mw^{[2]}+Bw$.
The quotient quadratic still has unrestricted rank and height $3h\le e+1$.

\subsection{Acquiring the ordinary phase at arbitrary rank}
The full binary centroid is exactly $A$, including after arbitrary higher
nilpotent cubic deformations. The residue form is radical-free: contraction
of $av+w$ has components $a\omega$ and $\ell\wedge\iota_w\omega$,
which vanish only for $a=w=0$. Its field centroid is scalar. In a balanced
map, repeated $v$ slots kill the component $v\to W$; three $W$ slots kill
the component $W\to v$; balancing $(v,w,w')$ forces its $W$ block to be
the same scalar as its $v$ block. The full-ring induction from
Section~\ref{sec:rankfive} gives the complete binary centroid $A$.
Raw acquisition therefore needs no supplied decomposition or bounded rank.

Freeze the complete acquired squarefree cubic
$C=\sum_{i<j<k}C_{ijk}X_iX_jX_k$. Its difference from the raw input is
binary quadratic. Fit the repeated-variable cubics $X_i^2X_j$, cross
pairs, linears and constant to every coefficient outside the low quotient.
For $K=b^2+\binom b2+b$, this system has $(K+1)ed$ columns and at most
$n(n+1)/2$ rows. Lemma~\ref{lem:invisible} applies in any number of
variables: simultaneous scaling isolates degree three, and an
odd-exponent variation isolates each coefficient. The complete protected
matrix has rank $Kd(e-h)$ and kernel dimension $ed+Kdh$.

For a fitted ordinary polynomial $P$, let $H=P_3$. Recover the center by
solving \emph{all} cross-pair equations at the necessary precision:
\[
 \sum_kC_{ijk}a_k=P_{X_iX_j}\pmod{J^{e-h}}.
\]
The residue contraction map has rank $b$, hence a unit minor and a left
inverse over every truncated ring. The center is unique modulo $J^{e-h}$.
Choose any lift, subtract $\lambda(H(X+a))$, and verify every remaining
binary coefficient and its quotient support, including affine terms.

Completeness follows by comparison with a promised representation.
Coefficient invisibility makes its fitted cubic and cross pairs agree
modulo $J^{e-h}$. Injectivity makes the centers agree to that precision.
Every term introduced by the remaining center difference contains a
factor from $J^{e-h}$ and is quotient-supported, even when the square
of that difference survives. Requiring the fitted cross vector to lie
in the full-ring contraction image would be wrong: for $b>3$, arbitrary
unconstrained deep cross coefficients need not lie in that image.
The program solves the stated congruence and checks all rows.

\subsection{Manufacturing the hyperplane and every candidate ray}
Solve $\alpha\wedge T_0=0$, a field-linear system with $b$ unknowns and
$\binom b4$ equations. Its kernel is exactly $F\ell$. Indeed, writing
$\alpha=a\ell+\beta$ on $W$ leaves $\beta\wedge\omega=0$.
For $\dim W\ge4$, this map is injective: a symplectic basis sends nonzero
$\beta$ to its first covector, whose product with the second symplectic
pair is nonzero. Conversely, $\ell\wedge T_0=0$ implies $T_0|_W=0$ and
$T_0=\ell\wedge\iota_vT_0$ for any $\ell(v)=1$. Check nonsingularity
on $W$. These linear systems acquire and certify the residue structure.

Lift a field basis of $W$ and a complement $v$ to the coefficient field
inside $A$. This remains an invertible ring coordinate change even when
the full higher-order third form has no contact factor. Write
\[
 H_0(w,z)=G_3(w)+zQ_2(w)+z^2L_1(w)+c_0z^3.
\]
Because $T_0|_W=0$, $G_3$ has only repeated-variable monomials and
$\nabla G_3=Mw^{[2]}$. The polar matrix of $Q_2$ is $B$.
Invertibility of $M$ excludes all nonzero critical vectors on $z=0$,
including over the algebraic closure. The single chart $z=1$ therefore
covers every rational projective critical point.

On that chart, $g(w)=Mw^{[2]}+Bw+c$ is an affine binary-linear system
on $(b-1)d$ bits. If consistent, stream its affine kernel coset of
$2^\kappa$ points and retain those with full field value $H_0(w,1)=0$.
Euler's cubic identity shows these are exactly the full critical rays.
Testing only the field trace would be wrong.

The bound $\kappa\le b-1$ is elementary. Kernel vectors satisfy
$w^{[2]}=Tw$, with $T=M^{-1}B$ invertible. A minimal field dependence
among binary-independent solutions, compared with its square and then
multiplied by $T^{-1}$, forces every dependence coefficient to satisfy
$a^2=a$. That contradicts binary independence. Thus binary-independent
kernel vectors are also field-independent.

Higher nilpotent mixed coefficients can make the full gradient nonlinear.
Its residue Jacobian is still the unit matrix $B$. At layer $t^j$, solve
$B\delta=[t^j]g$ and correct the point by $t^j\delta$. Verify every
full-ring transverse derivative after lifting. Retain the complete
critical value $c_v=H(v)$; only its residue is required to vanish.

\subsection{Counting with a bounded amount of working storage}
Let $\epsilon_v\in\{-1,1\}$ be the normalized Gaussian sign of the
formal quadratic Taylor part at a lifted ray, on its $2r$ transverse
field coordinates. The half-depth proof of Section~\ref{sec:criticalrays}
applies with an invertible $2r\times2r$ Hessian and gives
\[
 q^{re}\epsilon_v^{e\bmod2}
 \sum_{u\in A^\times}(-1)^{\lambda(c_vu^3)+Q(uv)}.
\]
For odd depth, scaling the quadratic by $u_0^3$ does not change its sign:
every nonzero field scalar has a square root, and scaling all transverse
variables by that root is a bijection. The weight is constant on the
surviving transverse cell because $h\le\lfloor e/2\rfloor$.

The radial phase is binary quadratic on $ed$ bits. Subtract its Gaussian
over $tA$ from its Gaussian over $A$; neither the rank of $Q$ nor the
module rank appears in an exponent here. The origin contributes
$q^{2b}Z_{e-3}$, with quotient height reduced by one. At depths two and
three it contributes $(-1)^{Q(0)}q^{b(e-1)}$.

The program builds the single origin path, then streams each binary-kernel
candidate and its Hensel lift across that path. It retains at most eight
sample rays and aggregate counts, not the exponential root list. There
are $O(e)$ origin states and two radial Gaussians per accepted ray per
state. All integers have $O(n)$ bits.

At a depth-one terminal, fix the acquired coordinate $z\in F$.
Every remaining cubic has a repeated variable, so $\Tr H_0(w,z)$ is
binary quadratic in $(b-1)d$ bits. Consequently
\[
 Z_1(H_0)=\sum_{z\in F}
   \operatorname{GaussianSum}_{w\in F^{b-1}}\Tr H_0(w,z).
\]
The implementation executes these $2^d$ Gaussian evaluations with an
explicit field-coordinate bit budget. It never enumerates $2^{bd}$ full
field assignments. For an input already at depth one, no ray enumeration
or infinity certificate is needed: the Gaussian fibers include $z=0$.
For fixed $b$, Lauder--Wan supplies another terminal option, but it is
not used to claim a bound uniform in $b$.

The acquired representation contains $n^2$ coordinate bits,
$\binom{b+2}{3}$ ring coefficients of $ed$ bits, and at most
$1+bdh+\binom{bdh}{2}$ quotient coefficients. All construction dimensions
are polynomial in $n+L$. One deliberately conservative bit bound is
\[
 O\!\left(Ln^2+n^{12}+
  \bigl(2^\kappa+\mathbf1_{e\equiv1\ (3)}2^d\bigr)n^{10}\right),
\]
with $O(L+n^6)$ working bits. Ring arithmetic, polynomial substitution,
binary elimination, finite-field operations and integer reconstruction
are included. For a pure field input the $2^\kappa$ term is unnecessary.
For a gate program, add the explicitly charged exact ANF propagation cost;
none of its wire values becomes an independent input.

\subsection{A polynomial family with unbounded module rank}
For any number of pairs $r\ge2$, use residue cubic
\[
 J_r=\sum_{i=1}^r
 (X_i^3+X_iY_i^2+Y_i^3+ZX_iY_i).
\]
Its square matrix has blocks
$\left(\begin{smallmatrix}1&1\\0&1\end{smallmatrix}\right)$ and its
Hessian has blocks
$\left(\begin{smallmatrix}0&1\\1&0\end{smallmatrix}\right)$.
The chart equations are
\[
 X_i^2+Y_i^2+Y_i=0,\qquad Y_i^2+X_i=0.
\]
They imply $Y_i(Y_i^3+Y_i+1)=0$. The cubic factor is irreducible over
$\F_2$, so for $3\nmid d$ there is exactly one solution, all transverse
coordinates zero. Thus $\kappa=0$ for every $b=2r+1$ and there is one
critical ray. This property survives arbitrary higher nilpotent
coefficients, arbitrary permitted quotient quadratics, and arbitrary
invertible binary affine disguises. It is acquired from raw input.

For $d\ge3$ with $3\nmid d$ and $e\equiv0,2\pmod3$, the complete
counter is polynomial uniformly in $b,d,e$. For all depths it is
polynomial when $d=O(\log n)$, using the executed Gaussian-fiber terminal.
There is no fixed module-width assumption. When $3\mid d$, each pair
instead contributes two binary kernel dimensions. The implementation
measures this change and charges it; it does not pretend the same
one-ray statement holds over those fields.

\subsection{A matching obstruction to ignoring the acquired kernel}
Arbitrary kernels make this geometric class \#P-hard even with fixed
$d=3$, $e=2$, nonsingular $M$, and no quotient noise. Let $F(u)$ be a
Boolean cubic on an even number $N$ of variables, padding by an unused
bit if needed. Put $m=N+2$, $b=m+1$, introduce two coordinates $a,\beta$,
and group $W=(u,a,\beta)$ into symplectic pairs. Set
\[
 q(W)=q_u(u)+a\beta,\qquad
 H_0=Zq(W)+\sum_{\text{pairs}}(X^2Y+XY^2).
\]
Its infinity matrix is nonsingular. The chart gradient forces every
coordinate of $W$ into $\F_2$, so $\kappa=m$. Let $G(W,Z)$ be the
ordinary degree-three homogenization of $F(u)+q_u(u)+a$, and put
$H=H_0+tG$. Every accepted critical ray has $q(W)=0$ and lifted
critical value $tG(W,1)$. At depth two its contribution is
$q^b(q-1)$ when $G=0$, and $-q^b$ when $G=1$, since cubing permutes
$\F_8^\times$.

For each fixed $u$, the two auxiliary coordinates give the exact identity
\[
 \sum_{a\beta=q_u(u)}(-1)^{q_u(u)+a}=1.
\]
The number of accepted rays is
$N_{\rm crit}=(2^m+2^{m/2})/2$. Including the origin yields
\[
 Z_A(H)=8^b\bigl(1+4Z_F+3N_{\rm crit}\bigr),
 \qquad n=6N+18.
\]
Thus one count recovers the arbitrary Boolean cubic signed sum by exact
division and subtraction. The full binary centroid is still $A$, and all
coefficients and original-input masks have polynomial description size.
The gate-constraint reduction in Section~\ref{sec:parityboundary} makes
this a \#P-hardness reduction. A polynomial counter uniform in arbitrary
$\kappa$ would imply $\mathrm{FP}=\#\mathrm P$.
This is not an unconditional $2^\kappa$ lower bound: some large-kernel
families have additional aggregations, including explicit formulas saved
with the audit sources.

There is also a separate obstruction if the infinity gate is dropped.
For $H_0=Z\sum_{i=1}^kX_iY_i$, $H=H_0+tG(X)$ and residue quadratic
$Q(X)$, direct highest-digit orthogonality gives
\[
 Z_A=q^{2k+1}\left(q^{k-1}S_f+
 (q^k-q^{k-1}+q-1)(-1)^{f(0)}\right),
\]
where $f(X)=\Tr G(X)+Q(X)$ and $S_f=\sum_X(-1)^{f(X)}$.
The positive-dimensional critical quadric retains an arbitrary cubic
count. The full reduction and its precise $24m+6$-bit encoding are
preserved in the independent audit.

\subsection{Executed raw counts, actual gates, and failures}
The raw contact compiler receives only the Boolean ANF and $n$.
The following accepted cases include centroid acquisition, the general-rank
ordinary fit, modular center recovery, hyperplane acquisition and exact counting:
\begin{center}
\small
\begin{tabular}{rrrrrrr}
\toprule
$n$&$b$&$d$&$e$&noise rank&$\kappa$&critical rays\\
\midrule
30&5&3&2&14&4&9\\
45&5&3&3&14&4&9\\
72&9&4&2&36&0&1\\
100&5&4&5&40&0&1\\
\bottomrule
\end{tabular}
\end{center}
Their signed sums are respectively $-2^{18}$, $-2^{30}$, $-2^{39}$,
and $-255\cdot2^{54}$. Every case passed $1024$ independent original-input
representation identities. The first two counts were checked by scanning
every residue vector and evaluating each complete surviving quadratic
fiber, without using projective rays or Hensel lifts. The 72-bit count
was checked by direct highest-digit orthogonality on its independently
proved critical line. The 100-bit case uses the undeformed one-ray
residue polynomial and an independent low-prefix enumeration oracle;
its actual raw recursion is $5\to2$.

A separate raw field case has $n=36$, $b=9$, $d=4$, $e=1$. Its signed
sum is $3997696$ and its zero count is $34361737216$. The new terminal
performs sixteen Gaussian eliminations on 32 variables, rather than
enumerating $2^{36}$ assignments. An independent pair-factor oracle
checks all $4096$ three-coordinate assignments; another $1024$ checks
verify the acquired raw representation. A field case with singular
infinity matrix is correctly accepted, because the complete Gaussian
fiber partition includes that hyperplane. An insufficient scan budget
is explicitly rejected.

Four small coefficient-level cases passed all $1098752$ assignments.
The independent geometry implementation passed eleven full truth tables,
totaling $4342784$ assignments, and verifies the exact change of kernel
dimension when $3\mid d$. Independent terminal tests covered six field
truth tables totaling $3195936$ assignments and two complete 20-bit
chain tables totaling $2097152$ assignments. These are correctness checks,
not a speed comparison or a claim of distinct assignments across suites.
The fixed-$\F_8$ hardness reduction was independently checked on all
$2129920$ residue assignments for 30- and 42-bit raw instances, with full
six-dimensional centroids. It does not claim their full raw truth tables.

The actual rank-nine gate description has 72 original inputs and 25678
provided gates. Exactly 8380 gates influence the output; backward
reachability is computed from the gate list. Exact propagation reaches
31476 ANF terms in one intermediate value and 31052 output terms.
Every output coefficient equals the independently generated raw input,
and 1024 original-input evaluations also match. The exact result is
\[
 Z=-549755813888,\qquad N_0=2361183241159944699904.
\]
The first frontend, which retained every wire polynomial, terminated
with exit code 137 and a recorded worker out-of-memory event. It produced
no count. The repaired frontend validates every gate, finds the output's
ancestors, derives all reference counts, and releases a polynomial after
its last use. Its peak logical cache contains 2706315 ANF terms on 279
wires. No original input or required computation is discarded.

For $G$ provided gates, maximum completed ANF size $K_g$, and peak live
term count $P_g$, a conservative deterministic bound for the literal
hash-set implementation is $O(L_c+GnK_g^4)$ bit time and
$O(L_c+n(P_g+K_g^2))$ working bits. This even charges worst-case hash
collisions and the temporary product parity set. Sorted products give
the usual smaller $O(GnK_g^2\log K_g)$ bound with the same quadratic
temporary-storage bound. The useful ANF sizes are measured, never assumed
small for all programs. Independent frontend checks covered 79 random
accepted programs, one rejected high-degree output, high-degree
cancellation, shared fanout and invalid dead gates, with 5376 original
assignments checked.

Two other failures are retained. A rank-seven nilpotent deformation
$T=Z\wedge\omega+tX_1\wedge X_2\wedge X_3$ has no full-ring contact
factor: a lift would require $X_1X_2X_3=\beta\wedge\omega$, whose
right side always contains a complete symplectic pair. The residue
method nevertheless lifts all 36 critical rays and retains their
nonzero higher critical values. Conversely, omitting the infinity
gate for $Z^3+Z(X_1Y_1+X_2Y_2)$ at $d=1,e=2$ would give 32 instead
of its exact count 320. The implemented gate prevents that loss.


\section{Manufacturing an eight-summary finite-field terminal}\label{sec:eightmoments}
The remaining field-coordinate scan can be removed from the complexity
theorem on a further acquired family. The number of module coordinates
is still unbounded. The extra certificate constructs a residue normal form
\begin{equation}\label{eq:pairnormal}
 H_0=\gamma\sum_{i=1}^r
 \left(X_i^3+X_iY_i^2+Y_i^3+
 Z(X_iY_i+a_iX_i^2+b_iY_i^2)\right)+\tau Z^3,
\end{equation}
with $\gamma\ne0$, $a_i,b_i\in\F_2$, and arbitrary $\tau\in F$.
The basis, scalar and pair types are manufactured. This includes $J_r$,
arbitrary invertible field coordinate changes and phase scaling, then
arbitrary higher nilpotent cubic coefficients, compatible quotient noise,
and the outer raw binary affine disguise. Require $d\ge3$ and $3\nmid d$.

\begin{theorem}[All-depth acquired finite-moment counter]\label{thm:eightmoments}
The preceding family has deterministic polynomial-time exact counting,
uniformly in $b=2r+1$, $d$, and $e$, with $n=bed$ original input bits.
A conservative bit-time bound is
\[
 O(Ln^2+n^{12})+\widetilde O(d^{31}),
\]
and working space is $O(L+n^6)+\widetilde O(d^{20})$.
The tilde terms occur only at a depth-one terminal. That terminal uses
seven explicit cubic point counts in at most eight field variables and
one binary Gaussian, followed by exact linear reconstruction.
\end{theorem}
The generic fixed-variable point-count engine is the established
Lauder--Wan algorithm \cite{lauderwan}; it is not implemented in this
source distribution. Acquisition, coefficient verification, the eight
reductions and reconstruction are implemented. The executed reference
backend explicitly scans a bounded field and is distinguished from the
polynomial-time theorem.

\subsection{Acquiring binary structure by semilinear descent}
Acquire the contact chart as before. Its affine gradient has a unique
binary solution for this family; solve it and change $w\mapsto w+az$.
The $z^2L_1$ term disappears. Write the centered polynomial as
$G_3(w)+zQ_2(w)+\tau z^3$, with square matrix $M$ and polar matrix
$\Omega$. Set $\sigma(w)=w^{[2]}$ and form
\[
 S=\Omega^{-1}M\sigma,\qquad K=S^d.
\]
The latter is field-linear and is computed by $d$ semilinear compositions.
Check $K^2+K+I=0$. Let $a\in\{1,2\}$ satisfy $ad\equiv1\pmod3$ and
put $D=K^a$, $F_*=D^{-1}S$. Acquire
$U=\ker_{\F_2}(F_*-I)$. Check that $\dim_{\F_2}U=b-1$ and that its
binary basis is a field basis of the transverse space.

In that basis, check that every entry of $\Omega$ is zero or one common
nonzero field scalar $\gamma$. The normalized form is binary symplectic.
Compute $D$ on $U$, check that its entries are binary and that it
preserves this form. These are polynomial-size identity tests, with no
enumeration of field elements or possible bases.

Manufacture $D$-invariant symplectic planes. Since $D^2+D+I=0$, the
binary quadratic $v\mapsto\Omega(v,Dv)/\gamma$ has polar form
$\Omega/\gamma$. It is nonzero, so a vector with value one is found
among a basis and its pairwise sums. The plane spanned by $v,Dv$ is
nonsingular, and its orthogonal complement is $D$-invariant. Repeat.
Every acquired plane has blocks
\[
 \Omega/\gamma=\begin{pmatrix}0&1\\1&0\end{pmatrix},\qquad
 D=\begin{pmatrix}0&1\\1&1\end{pmatrix}.
\]
Because the basis lies in $U$, $S=D\sigma$ there, hence $M=\Omega D$.
This gives exactly the repeated cubic blocks in \eqref{eq:pairnormal}.
Read the square coefficients of $Q_2$ and require their quotients by
$\gamma$ to be binary. Finally verify the \emph{entire ordinary
polynomial identity}, not a sample of Boolean values.

To prove completeness under hidden coordinates, start in the displayed
family. After acquiring the contact line and transverse center, the old
and new coordinates are $w_{\rm old}=Rw$, $z_{\rm old}=tz$, with $R$
invertible and $t\ne0$. If $D_0$ is the sum of the displayed order-three
blocks and the original scalar is $\gamma_0$, then
\[
 S=t^{-1}R^{-1}D_0\sigma R,\quad
 K=R^{-1}D_0^dR,\quad D=R^{-1}D_0R,\quad
 F_*=t^{-1}R^{-1}\sigma R.
\]
The scalar norm of $t$ is one. Therefore
$U=R^{-1}t\F_2^{b-1}$ and the common symplectic scale is
$\gamma_0t^3$. Every test succeeds. The acquired planes may differ
from the hidden ones, but the normalized square coefficients remain
binary and the full coefficient identity remains valid. The earlier
centroid and quotient-fit proof handles the outer binary affine disguise.

The rotation $(X,Y)\mapsto(Y,X+Y)$ preserves the cubic block and sends
square type $(a,b)$ to $(b,1+a+b)$. It cycles $(0,0),(0,1),(1,0)$ and
fixes $(1,1)$. Thus only two counting types remain, with acquired
multiplicities $r_+$ and $r_-$.

\subsection{Symmetry gives eight joint Gaussian states}
For fixed $z\in F$, define
\[
 A_\pm(z)=\sum_{x,y\in F}(-1)^{\Tr\{\gamma[
 x^3+xy^2+y^3+z(xy+\delta_\pm(x^2+y^2))]\}},
 \quad\delta_+=0,\quad\delta_-=1.
\]
The phases are binary quadratic on $2d$ bits with the same polar matrix.
Squaring their common polar-kernel equations gives
\begin{align*}
 \gamma x^4+x+\gamma y^4+\gamma z^2y^2&=0,\\
 x+\gamma y^4+y+\gamma z^2x^2&=0.
\end{align*}
The ordinary Jacobian is
$\left(\begin{smallmatrix}1&0\\1&1\end{smallmatrix}\right)$.
All geometric solutions are isolated and simple; the degree bound gives
at most sixteen. The common binary polar nullity $k$ is therefore one
of $0,2,4$. Each Gaussian is zero or $\pm q2^{k/2}$.

The minus phase is invariant under $C(x,y)=(y,x+y)$, with
$C^2+C+I=0$. Its radical is $C$-stable. The quadratic is linear there;
invariance and $v+Cv+C^2v=0$ force it to vanish on the radical.
On the nonsingular quotient, manufacture $C$-invariant symplectic planes.
In each plane spanned by $v,Cv$ with polar pairing one, invariance forces
the quadratic to equal one on all three nonzero vectors. Each plane has
negative Gaussian sign. Consequently
\[
 A_-(z)=(-1)^{d-k/2}q2^{k/2}\ne0.
\]
When $k=0$, $A_+$ is also nonzero. With $\varepsilon=(-1)^d$, the joint
state $(A_+,A_-)$ belongs to the eight-element set
\[
 \mathcal V=\{(\pm q,\varepsilon q),\ (0,-2\varepsilon q),
 (\pm2q,-2\varepsilon q),\ (0,4\varepsilon q),
 (\pm4q,4\varepsilon q)\}.
\]
Its size is independent of $b,d,e$; each integer entry has $O(d)$ bits.

\subsection{Eight manufactured summaries and exact reconstruction}
For $(u,v)\in\mathcal V$, define the signed multiplicity
\[
 C_{u,v}=\sum_{z:(A_+(z),A_-(z))=(u,v)}(-1)^{\Tr(\tau z^3)}.
\]
Use moments $M_{ij}=\sum_z(-1)^{\Tr(\tau z^3)}A_+(z)^iA_-(z)^j$
at the eight indices
\[
 (0,0),(0,1),(0,2),(1,0),(1,1),(1,2),(2,0),(2,1).
\]
Their evaluation matrix on $\mathcal V$ is invertible over the rationals.
Group the nodes by their three distinct second coordinates. The first
three moments recover the groups' signed masses; the next three recover
their first $u$-moments. In the first group $u^2=q^2$ identically.
Subtract its known contribution from the last two moments and solve the
two-by-two system for the remaining second $u$-moments. These determine
each zero-node mass and the separate masses at opposite nonzero values.
This proves uniqueness and supplies exact reconstruction. The determinant
magnitude of the eight-by-eight matrix is $2^{15}3^7q^{14}$.

The field count is then
\[
 Z_1=\sum_{(u,v)\in\mathcal V}C_{u,v}u^{r_+}v^{r_-},
\]
with $0^0=1$. Moment and reconstruction operands have $O(d)$ bits.
The final powers and counts have $O(bd)$ bits. No unit-cost operation
on an exponentially long integer is assumed.

Each positive-degree moment expands into at most three pair copies,
sharing $z$. It is the character sum of the explicit ordinary cubic
\[
 P_{ij}=\tau z^3+\gamma\sum_{\ell=1}^{i+j}
 [x_\ell^3+x_\ell y_\ell^2+y_\ell^3+
 z(x_\ell y_\ell+\delta_\ell(x_\ell^2+y_\ell^2))],
\]
where the first $i$ copies have $\delta_\ell=0$ and the last $j$ have
$\delta_\ell=1$. It has $2(i+j)+1\le7$ field variables. Manufacture
the Artin--Schreier equation $w^2+w+P_{ij}=0$, giving
\[
 M_{ij}=\#\{w^2+w+P_{ij}=0\}(F)-q^{2(i+j)+1}.
\]
These seven hypersurfaces have degree three and at most eight variables,
independently of the number of pairs. The zeroth moment is the binary
Gaussian $\sum_z(-1)^{\Tr(\tau z^3)}$.
Lauder--Wan's bound for fixed characteristic, degree and variable count
$s$ is $\widetilde O(d^{3s+7})$ time and
$\widetilde O(d^{2s+4})$ space \cite{lauderwan}. Taking $s\le8$ gives
the stated $d^{31}$ and $d^{20}$ bounds. The summaries are therefore
manufactured by fixed-dimensional computations, not assumed as advice.

The residue gradient in \eqref{eq:pairnormal} has the same zero binary
kernel as $J_r$ when $3\nmid d$. There is one projective critical ray
if $\tau=0$, and none if $\tau\ne0$. All ring stages are polynomial;
this terminal closes the remaining depth-one case and proves
Theorem~\ref{thm:eightmoments}.

\subsection{A sharp divisibility boundary in the same residue family}
The restriction $3\nmid d$ cannot simply be deleted while keeping arbitrary
higher nilpotent cubics. The very same $J_r$ residue already gives a
\#P-hard family at the fixed degree $d=3$ and depth $e=2$, with no quotient
quadratic. This is stronger than hardness in the broader contact class.

Let $p$ be a Boolean cubic on an even number $N\ge2$ of variables, and set
$r=N+1$, $b=2r+1=2N+3$. In $F=\F_8$, choose
$\alpha^3+\alpha+1=0$. A nonzero critical residue of $J_r$ has $Z\ne0$;
normalize $Z=1$. Its transverse equations give
\[
 X_i=Y_i^2,\qquad Y_i^4+Y_i^2+Y_i=0,
 \qquad Y_i=a_i\alpha+b_i\alpha^2,\quad a_i,b_i\in\F_2.
\]
The remaining derivative is $\sum_iX_iY_i=\sum_iY_i^3=0$.
The exact identity
\[
 Y_i^3=a_i\alpha+b_i\alpha^2+(a_i+b_i+a_ib_i)
\]
therefore makes the critical-ray constraints
\[
 \sum_i a_i=0,\qquad \sum_i b_i=0,\qquad \sum_i a_ib_i=0.
\]
The field-linear form $L_i=\alpha^2X_i+\alpha^4Y_i$ equals $a_i$ on
this chart. Manufacture the ordinary homogeneous cubic
\[
 G(X,Y,Z)=\sum_{S\in\operatorname{supp}(p)}
 Z^{3-|S|}\prod_{i\in S}L_i,
 \qquad H=J_r+tG\in F[t]/(t^2)[X,Y,Z].
\]
It has only polynomially many explicitly constructed coefficients.
On a critical chart point, $G=p(a_1,\ldots,a_N)\in\F_2$.

Because $r$ is odd, an even-parity vector $a$ cannot be the all-ones
vector. For $a\ne0$ the two equations on $b$ are independent and have
$2^{r-2}$ solutions; for $a=0$ there are $2^{r-1}$. The first $N$
coordinates determine the last $a$ uniquely. If
$Z_p=\sum_x(-1)^{p(x)}$ and $s_0=(-1)^{p(0)}$, then
\[
 N_{\rm crit}=2^{N-1}(2^N+1),\qquad
 W_{\rm crit}=2^{N-1}(Z_p+s_0).
\]
Summing all highest ring digits gives the factor $8^b$ and restricts the
residue to $\nabla J_r=0$. The origin contributes one. On each nonzero
ray, cubing permutes $F^\times$, so the radial sum is seven when $G=0$
and minus one when $G=1$, namely $3+4(-1)^G$. Consequently
\begin{equation}\label{eq:samepairhardness}
 Z(H)=8^{2N+3}
 \left[1+3\,2^{N-1}(2^N+1)
       +4\,2^{N-1}(Z_p+s_0)\right],
 \qquad n=12N+18.
\end{equation}
One exact count recovers $Z_p$ by polynomial-bit arithmetic. The full
binary centroid remains $F[t]/(t^2)$ by the earlier contact-centroid
proof, and arbitrary invertible binary affine disguises preserve the
count. Thus this is a raw acquired-algebra obstruction, not a supplied
decomposition counterexample.

For completeness, signed Boolean cubic counting is \#P-hard directly:
encode each bounded-fan-in Boolean gate by a quadratic equation
$g_j(x,y)=0$, include the accepted-output equation, and introduce a
fresh binary multiplier $u_j$. The phase $\sum_j u_jg_j$ is cubic.
Summing the multipliers enforces every equation and multiplies the
number of accepted computations by $2^{\#j}$. Gate values are uniquely
determined by the input. Add one unused bit if necessary to make $N$
even and divide its known factor out. This proves the asserted reduction
without assuming a lower bound on every large-kernel instance.

The executable manufactures two raw affine disguises. The 42-bit case
has ten critical rays, source signed count $-2$ and ring signed count
$14680064$; it independently scans all $2097152$ residue vectors.
The 66-bit case has 136 rays, source signed count zero and ring signed
count $3238405341184$; it independently scans 1024 chart candidates.
Both recover the source count through \eqref{eq:samepairhardness},
acquire a six-dimensional binary centroid, and pass 256 original-input
representation checks. Neither test enumerates its full ring domain.
An independent coefficient audit also covers the 90-bit construction.

\subsection{Execution boundary and prior machinery}
The source implements semilinear descent, plane acquisition, every
ordinary coefficient check, the eight-state reconstruction, and explicit
Artin--Schreier polynomial manufacture. Six disguised coefficient-field
cases passed $2114720$ complete assignments; twenty independent expanded
moments covered $6624$ assignments. Five larger acquisition-only cases
reach 77 binary residue coordinates, with no field scan during acquisition.
The raw 36-bit case reacquires the ring and phase, then the pair certificate
with $\gamma=10$, $\tau=0$, and multiplicities $(4,0)$. Its signed count
is again $3997696$.

That raw test obtains reference moments from an explicitly recorded
sixteen-element field scan and replays seven through the moment interface;
the zeroth is a Gaussian. A further raw 80-bit case at $d=4,e=4,b=5$
includes arbitrary higher cubic coefficients and rank-20 quotient noise.
It reacquires the pair template after its actual depth reduction and
agrees with an independently exhausted field leaf. An additional audit
checks 406 pair-Gaussian cases, covering 327056 phase evaluations and
every polar nullity $0,2,4$.

The production hypersurface adapter has also been checked by evaluating
112 actual equations over $\F_2$ and $\F_4$, including their
Artin--Schreier coordinates: $1728136$ complete assignments. The resulting
cardinalities agree with $432404$ separately expanded character evaluations.
The raw 36-bit callback test receives exactly the seven manufactured
coefficient dictionaries, with at most eight variables, and reproduces
the raw count. That callback transparently replays moments prepared by
the recorded sixteen-element field scan. For an opaque supplied backend,
the wrapper reports only its own enumeration count, not the backend's
unobserved internal work.

These executions validate acquisition, reductions and reconstruction.
They do not execute the generic Lauder--Wan engine. The raw wrapper
accepts a point counter receiving each manufactured hypersurface, or
labels its bounded reference scan. The earlier 49-moment construction
is preserved as an independently checked reference; the eight-moment
method is the production interface.

Quadratic-rank and finite-moment methods have established ancestry in
coding theory, including Luo \cite{luomoments}. For odd $d$, writing
$u=x+\omega y$ in $\F_{q^2}$ with $\omega^2+\omega+1=0$ identifies
these pair sums with a generalized Kasami sum and its linked linear
twist studied by Luo--Tang--Wang \cite{kasami}. The source proof gives
the exact identification. No invention of those classical sums or of
moment inversion is claimed. The candidate contribution is the
manufactured semilinear and block representation, the eight summaries
whose dimension is independent of unbounded module rank, and their
proved composition with raw weighted chain-ring counting. Historical
priority remains unestablished.


\section{Manufacturing a global defect support}\label{sec:sylvester}
Write $V=F^b$, where $b=3e$ for the local-algebra compiler, and use
$\langle x,y\rangle=\Tr\sum_i x_i y_i$. Scalar multiplication $M_a$ is
selfadjoint for this pairing. A quadratic polar matrix gives
$L=G^{-1}B$. The compatible bandwidth-one space consists of
\[
 L_0(x)=C_{-1}x^{2^{d-1}}+C_0x+C_1x^2,\qquad C_j\in F^{b\times b}.
\]
Every ordinary ring cubic, expanded in the coefficient field, has its
quadratic polar part in this space. Repeated variables give Frobenius
powers $0,\pm1$; the nilpotent convolution changes only coefficient matrices.

Let $z$ be the acquired degree-$d$ field generator and set
\[
 c=z+z^2,\quad v=z+z^3,\quad
 C=\Span_{\F_2}\{1,z^4,c,v\},\quad S=C+cC.
\]
Consider the symmetric bivariate polynomial
\begin{align*}
 P(X,Y)&=(X+Y)(X^2+Y)(X+Y^2)\\
 &=X^4+Y^4+(X+X^2)(Y+Y^3)+(X+X^3)(Y+Y^2).
\end{align*}
Manufacture the operator
\begin{equation}\label{eq:selfadjoint}
 H(L)=LM_{z^4}+M_{z^4}L+M_cLM_v+M_vLM_c.
\end{equation}
It is computed from $B$ and the acquired field.

\begin{lemma}
For every $d\ge7$, $\ker H$ is exactly the bandwidth-one space,
$\dim C=4$, $C\cap cC=\F_2c$, and $\dim S=7$.
\end{lemma}
\begin{proof}
Every binary linear map has a unique matrix-coefficient linearized-polynomial
expansion $\sum_t C_t x^{2^t}$. On its $t$th term, $H$ multiplies the
coefficient by $P(z,z^{2^t})$. The degree-$d$ Frobenius orbit has distinct
elements; exactly $t=0,1,-1$ are zeros of the three displayed factors.

For the scalar claims use polynomials in an indeterminate $Z$. The four
basis elements have distinct highest degrees $0,4,2,3$. Of their $c$
multiples, $c$ is already in $C$; $c^2=Z^2+Z^4$ is outside $C$;
$cv$ has degree five and $cZ^4$ degree six. These last three are
independent modulo $C$. No nonzero polynomial of degree at most six
vanishes at $z$ when $d\ge7$.
\end{proof}

\begin{theorem}[Global support recovery]
Suppose $L=L_0+E$, where $L_0$ is compatible, $GE$ is alternating of rank
$r$, and $\dim(S\im E)=7r$. Then $L_0$ and $E$ are recovered exactly from
$L$ in deterministic polynomial time. No inequality between $r$ and $d$
is needed.
\end{theorem}
\begin{proof}
For a binary basis matrix $U$ of $\im E$, selfadjointness gives
$E=UKU^*$ with $K$ nonsingular alternating and $U^*=U^\top G$.
Set $V_C=[U\ \ z^4U\ \ cU\ \ vU]$, and let $J$ swap its first two
and last two block positions. Then
\[
 H(L)=H(E)=V_C(J\otimes K)V_C^*.
\]
Directness on $S\otimes U$ makes $V_C$ full column rank, so
$W=\im H(L)=CU$. It also preserves intersections:
\[
 W\cap cW=(C\cap cC)U=cU,\qquad
 U=c^{-1}(W\cap cW).
\]
This manufactures the support. Express every column of $H(L)$ uniquely
in $[U\ z^4U\ cU\ vU]$. Its second block of coefficient rows is
$KU^*$; multiplication by $U$ recovers $E$. All steps are binary
linear algebra and known scalar multiplication.
\end{proof}

The executable checks the recovered support's sevenfold directness,
alternation and rank of $GE$, and that the reconstructed base lies in the
Frobenius band. These checks certify every returned decomposition, even
outside the sufficient planted-recovery hypothesis. A zero syndrome returns
a checked rank-zero alternative; it cannot identify an invisible error
that was itself compatible.

For independent uniform columns $u_1,\ldots,u_r\in F^b$, any nonzero
$(s_1,\ldots,s_r)\in S^r$ makes $\sum_i s_i u_i$ uniform on $F^b$:
condition on all but one column with nonzero scalar coefficient.
Thus a union bound gives
\[
 \Pr[\dim(SU)<7r]\le(2^{7r}-1)2^{-n}<2^{7r-n}.
\]
Conditioning $U$ to have rank $r\le n-1$ costs at most a factor two.
With a fixed nonsingular alternating $K$, uniform full-rank $U$ induces
the uniform distribution on rank-$r$ alternating matrices after the fixed
Gram transformations, by transitivity of binary congruence and equal fiber
sizes. The acquired coordinates depend only on the third tensor, not the
quadratic error, so the same bound applies before the raw disguise is undone.
No random-subspace intersection heuristic is used.

The preceding one-sided construction used
$Q(T)=\prod_{j=0}^2(T+z^{2^j})$ and
$\sum_jM_{[T^j]Q}LM_{z^2}^j$. It required a ten-dimensional combined scalar
orbit and uniformly $d\ge12$. Its proof, executed successes and failures
are retained. Selfadjointness removes the separate right orbit, yielding
the stronger seven-direction theorem. For larger bandwidths the preserved
one-sided construction remains available.

\section{Using nilpotence itself to recover the defect}
Assume $e\ge7$, and use the ring pairing
$\langle x,y\rangle_A=\lambda(\sum_i x_i y_i)$ with Gram matrix $G_A$.
Let $N$ be multiplication by the acquired uniformizer. It is selfadjoint.
Replace $z$ by $N$ in \eqref{eq:selfadjoint}, and put
\[
 C=\Span\{I,N^4,N+N^2,N+N^3\},\quad c=N+N^2,\quad
 S=C+cC=\Span\{I,N,\ldots,N^6\}.
\]
The degree-six coefficient argument gives $\dim S=7$ and
$C\cap cC=\F_2c$.
If $Q$ denotes squaring on $A$, then $QN=N^2Q$.
Thus the operator annihilates $M_aQ$ and, by taking adjoints, $Q^*M_a$.
It also annihilates multiplication maps. These are the types occurring
in ordinary ring-cubic polars. Its kernel may be larger than that base
family: exact ordinary-cubic fitting and factor verification remain essential.

For $E=G_A^{-1}D$ and $U=\im E$, assume $\dim(SU)=7r$.
The same selfadjoint factorization gives $W=CU$.
Although $c$ is not globally invertible, it is injective on $CU$:
the four directions of $cC$ are independent within $S$.
Consequently
\[
 U=\{w\in W:cw\in W\}.
\]
Indeed, $cw\in W\cap cW=cU$ implies $cw=cu$ for some $u\in U$,
and injectivity on $CU$ implies $w=u$. A restricted linear solve
computes the preimage. The second syndrome-coordinate block then recovers $E$.

\begin{corollary}[Long-chain counter over small residue fields]
For every $d\ge3$, $e\ge7$, the raw family with a rank-$r$ defect is counted
in $2^r\poly(L+n)$ time and polynomial space when the acquired uniformizer
support satisfies $\dim(SU)=7r$. For iid uniform support columns the failure
probability is at most
\begin{equation}\label{eq:nilprob}
 \sum_{v=0}^6
 \left(2^{(7-v)r}-2^{(6-v)r}\right)2^{-n+3dv}.
\end{equation}
For uniform rank-$r$ alternating defects multiply by at most two.
\end{corollary}
\begin{proof}
Classify nonzero tuples $(s_1,\ldots,s_r)\in S^r$ by minimum valuation $v$.
There are $2^{(7-v)r}-2^{(6-v)r}$ such tuples.
One coefficient is $N^v$ times a unit, so conditional on all other columns
the relation is uniform on $N^vV$, of dimension $n-3dv$.
A union bound gives \eqref{eq:nilprob}. Exact recovery and evaluation
follow the algorithms above and below.
\end{proof}
A simpler bound is $2^{7r-n+18d}$ before conditioning. Fixed $d$,
$e\to\infty$ and $r=o(n)$ give success probability $1-2^{-n+o(n)}$.
The field and nilpotent conditions refer to the generator actually acquired.
A generator change need not preserve a particular scalar space, but
acquisition's independence of the error proves the probability claims.

\section{Recovering the ordinary ring cubic and its factors}
After error removal, subtract the canonical alternating cubic
$\lambda(\mu XYZ)$. Fit the remaining quadratic polar coefficients by a
binary linear system. Candidate ring monomials are $XY,XZ,YZ$,
$X_i^3$, and $X_i^2X_j$ for $i\ne j$. Let each coefficient range over
the acquired binary ring basis. There are $12ed=4n$ scalar columns and
$\binom n2$ polar equations. Unary and constant corrections are then solved
using the nondegenerate ring pairing. A successful solve is an exact
Boolean polynomial identity, not agreement on sampled points.

\begin{lemma}[Identifiability of the homogeneous cubic]
For $d\ge3$, a homogeneous ordinary cubic $H\in A[X_1,\ldots,X_b]$ is
uniquely determined by its scalar function $\lambda(H)$ modulo scalar
traces of ordinary polynomials of degree at most two.
\end{lemma}
\begin{proof}
Fix a ring point $X$ and scale all its coordinates by $\rho\in F$.
The Frobenius orbit of exponent three modulo $2^d-1$ has length $d$
and is disjoint from the orbits of zero, one and two. Therefore
$\Tr(\rho^3 a)$ cannot cancel those lower-degree trace functions unless
$a=0$. This follows either from reduced polynomial coefficients or from
the distinct cyclotomic exponent classes. Consequently
$[t^{e-1}]H(X)=0$ for every $X$ if the scalar cubic is a lower-degree trace.

Expand each $X_i=\sum_j t^j x_{ij}$. The top coefficient is an ordinary
polynomial of total degree three in the field coefficient variables;
degree three is below $q=2^d$, so zero evaluation means the polynomial is
identically zero. For a coefficient $h_{\alpha,j}$ of $t^jX^\alpha$,
choose $i$ with $\alpha_i$ odd and put $k=e-1-j$. If $k>0$, the component
monomial
$x_{i,k}x_{i,0}^{\alpha_i-1}\prod_{\ell\ne i}x_{\ell,0}^{\alpha_\ell}$
has coefficient exactly $h_{\alpha,j}$. No other exponent or ring layer
contributes. If $k=0$, use $x_0^\alpha$. Hence every coefficient vanishes.
\end{proof}

An acquired ring isomorphism or free module basis change preserves ordinary
degree. The original Frobenius functional becomes $\lambda(c\,\cdot)$ for
a unit $c$, absorbed into coefficients. Thus the fit recovers the hidden
homogeneous cubic despite arbitrary binary disguise.

Reduce the homogeneous cubic modulo $t$. Its three independent field-linear
factors are acquired using the preserved deterministic field routine:
two seven-probe line searches produce a squarefree cubic slice, trace-gcd
splitting finds its roots in polynomial time in $d$, and gradients give
the factor normals. The routine verifies every polynomial coefficient.
No search through all $2^d$ field elements is required.

Send these residue factors to $X,Y,Z$. Write the sought lifted factors as
three linear forms with diagonal coefficients one and off-diagonal
coefficients in $tA$, multiplied by a unit coefficient $\mu$.
At each layer $t^j$, the correction has the form
\[
 \delta\mu\,XYZ+\mu_0\bigl(\delta L_1\,YZ+
              X\,\delta L_2\,Z+XY\,\delta L_3\bigr).
\]
The six off-diagonal coefficients and $\delta\mu$ independently control the
six repeated-variable mixed cubics and $XYZ$. The three pure-cube
coefficients are obstruction tests. Solve these seven coefficients at each
of the $e-1$ layers and verify the final full factor identity.
Independence of the residue planes makes the lifted matrix invertible.
Under the split-family promise, the normalization only fixes harmless
unit-scaling gauges, so the lifts exist at every stage.

Ring squares must be retained before counting. If
$c=\sum_j c_jt^j$, the correct scalar adjoint of squaring is
\begin{equation}\label{eq:square}
 \lambda(cX^2)=\lambda(c^\dagger X),\qquad
 c^\dagger=\sum_{2j\le e-1}\sqrt{c_{e-1-2j}}\,t^{e-1-j}.
\end{equation}
This follows by tracing the top coefficient term by term. Reusing the
field-only formula $c^{2^{d-1}}$ is wrong when nilpotents are present.
The full proof and an executed regression are preserved.

\section{An actual polynomial local-ring evaluator}
Write $q=2^d$, $R_e=F[t]/t^e$, and
$\psi_e(x)=(-1)^{\Tr([t^{e-1}]x)}$.
We evaluate
\[
 \sum_{x,y,z\in R_e}\psi_e(\mu xyz+Axy+Bxz+Cyz+Dx+Ey+F_1z+G_0)
\]
for every unit $\mu$ and arbitrary other coefficients.
First scale $z'=\mu z$ and replace $B,C,F_1$ by their quotients by $\mu$.
Then translate $x=X+C$, $y=Y+B$, $z'=Z+A$. The polynomial becomes
$XYZ+aX+bY+cZ+g$, where
\[
 a=D+AB,\quad b=E+AC,\quad c=F_1+BC,\quad
 g=G_0+DC+EB+F_1A.
\]
All substitutions are bijections. Summing over $Z$ yields
\begin{equation}\label{eq:productfiber}
 Z=q^e\psi_e(g)S_e(a,b;c),\qquad
 S_e(a,b;c)=\sum_{xy=c}\psi_e(ax+by).
\end{equation}

Put $v(0)=e$, and
$K_e(A,D)=\sum_{x\in R_e^\times}\psi_e(Ax+D/x)$.
Decompose $x=t^ku$ by valuation. The fibers in $y$ have $q^k$ elements;
their character sum is zero unless $v(b)\ge k$. Hence
\begin{equation}\label{eq:valuation}
 S_e(a,b;c)=q^e\,1_{b=c=0}+
 \sum_{k=0}^{\min\{v(b),v(c),e-1\}}
 q^k K_{e-k}\left(a\bmod t^{e-k},
            (b/t^k)(c/t^k)\bmod t^{e-k}\right).
\end{equation}
The first term includes $x=0$. This also gives
$q^e+e(q-1)q^{e-1}$ pairs satisfying $xy=0$.

For $s=\min(v(A),v(D))<e$,
$K_e(A,D)=q^sK_{e-s}(A/t^s,D/t^s)$.
If both coefficients vanish the sum is $(q-1)q^{e-1}$.
After reduction, exactly one unit coefficient gives $-1$ at length one
and zero at length at least two, by summing the highest coefficient of $x$.
Two unit coefficients reduce to $K_e(1,AD)$ by scaling $x$.
At length one the executable uses the actual finite Dwork evaluator from
the preserved field compiler; it does not enumerate $F$.

\subsection{Manufacture the stationary center}
For a unit $u$ and $e\ge2$, put
$f(x)=x+u/x$, $c=\lceil e/3\rceil$ and $z=\lceil(e-1)/3\rceil$.
For $y\in t^cR_e$,
\[
 f(a+y)-f(a)=U(a)y+V(a)y^2,\qquad U(a)=1+u/a^2,\quad V(a)=u/a^3,
\]
because $y^3=0$. The character on this coset is constant precisely when
for $0\le i\le e-1-c$,
\begin{equation}\label{eq:center}
 U(a)_i+1_{2i-e+1\ge0}\sqrt{V(a)_{2i-e+1}}=0.
\end{equation}
Indeed, the coefficient $y_j$ contributes
$U_{e-1-j}y_j+V_{e-1-2j}y_j^2$ to the top trace, and
$\Tr(vy^2)=\Tr(\sqrt v\,y)$.

These conditions depend only on $a\bmod t^z$. In the mixed second
difference for shifts in $t^c$ and $t^z$, the first possible cross terms
have valuations $2c+z$ and $c+2z$, both at least $e$.
The zeroth equation fixes $a_0=\sqrt{u_0}\ne0$. To find $a_j$ for
$1\le j<z$, temporarily set it to zero and use equation $i=2j$.
The needed index $4j-e+1<j$ in $V$ is already determined. Replacing
$a_j$ by $\alpha$ changes $U_{2j}$ by $\alpha^2/a_0^2$ and no earlier
equation. Thus the only possible coefficient is
\[
 a_j=a_0\sqrt{U(a)_{2j}
             +1_{4j-e+1\ge0}\sqrt{V(a)_{4j-e+1}}}.
\]
All these indices lie in the tested range. Construct the coefficients
and then verify every equation \eqref{eq:center}, including odd indices.
Any failed equation makes $K_e(1,u)=0$.

If $c=z$, exactly one coset survives and
$K_e(1,u)=q^{e-c}\psi_e(f(a))$.
Otherwise $e=3z+1$ and $c=z+1$. The surviving cosets have representatives
$a+t^zw$ for $w\in F$, giving
\begin{equation}\label{eq:gauss}
 K_e(1,u)=q^{e-c}\psi_e(f(a))
       \sum_{w\in F}(-1)^{\Tr(\gamma w^3+\beta w)},
\end{equation}
where $\gamma=a_0^{-2}$ and $\beta=U(a)_{2z}+\sqrt{V(a)_z}$.
Although its ordinary degree is three, $w^3=w\,w^2$ is quadratic over
$\F_2$. Evaluate at zero, binary basis points and pairs to acquire its
quadratic form, then eliminate quadratic pairs exactly. This computes
the sign, radical and zero cases without enumerating the field.

\subsection{Arithmetic cost and a source-bound failure}
A ring inverse uses $O(e^2+d)$ field operations by coefficient recurrence.
The center uses $O(e)$ such steps, and the terminal quadratic costs
$O(d^3)$ field operations and $O(d^4)$ bit work. Formula
\eqref{eq:valuation} makes at most $e$ Kloosterman calls.
With schoolbook field arithmetic and the preserved $O(d^8)$ bit
Dwork backend, a conservative total is
\[
 O(e^4d^2+ed^8)+\poly(ed)
\]
bit operations and polynomial space. Exact answers have $O(ed)$ bits.
These are charged arithmetic bounds, not timing claims.

Local stationary phase has established ancestry, particularly Sawin's
work on wild Kloosterman sums \cite{sawin}. We do not use its displayed
magnitude bound: the inspected arXiv v2 has a characteristic-two boundary
counterexample. In $\F_2[t]/t^4$, set $u=1+t^2$.
For $x=1+at+bt^2+ct^3$, the top coefficient of $x+u/x$ is
$a^3+a=0$, so all eight unit terms are $+1$ and $K_4(1,u)=8$.
The displayed Corollary 1.4 specializes to $2^{5/2}$.
Our counter evaluates the nontrivial terminal quadratic radical exactly.
The full eight-term computation is preserved. No erratum was located in
the targeted search; this does not establish that no correction exists.

\section{Reconstruction, full costs, and actual programs}
Exact quadratic elimination writes the canonical error as
\[
 q_D(y)=\sum_{j=1}^{r/2}\ell_{2j-1}(y)\ell_{2j}(y)+\ell_0(y).
\]
The factors are manufactured, and all unary corrections are retained.
The elementary identity
\[
 (-1)^{uv}=\frac12\sum_{\alpha,\beta\in\F_2}
     (-1)^{\alpha\beta+\alpha u+\beta v}
\]
gives $2^r$ signed linear-character leaves and denominator $2^{r/2}$.
Every binary parity in acquired ring coordinates is uniquely
$\lambda(a_1X+a_2Y+a_3Z)$. After the acquired factor substitution it is
still ring-linear. Hence each leaf is an instance of the exact evaluator
above, with its original cubic factors unchanged.

Accumulate the signed numerator using $O(n+r)$ bits and divide exactly
only after summing. Leaves are streamed, not stored. Convert to
$N_0=(2^n+Z)/2$. A convenient loose full bit bound, using explicit dense
binary elimination, is
\begin{equation}\label{eq:cost}
 O\!\left(n^7+Ln^2+
    2^r\bigl(n^4+e^4d^2+ed^8\bigr)\right),
\end{equation}
with $O(L+n^4+d^4)$ working bits up to standard representation overhead.
The $n^7$ allowance charges the centroid equations and all algebra,
support, coefficient and factor computations; no algebraic operation is
declared free. The public representation itself needs an $n\times n$ binary
coordinate matrix, ring parameters, constant-size ring polynomial,
and rank-$r$ support data, using $O(n^2+nr+ed)$ bits.

The default compiler tries the applicable Gabidulin, selfadjoint-field and
nilpotent-support routes. Each candidate must pass the ordinary-cubic fit
and full factor identity. Among verified candidates it selects the smallest
defect rank. There are at most three polynomial-cost trials. This acquisition
choice is executed; a user need not provide the useful rank or method.

\subsection{Primitive products}
The fixed space of squaring contains the idempotents of a commutative
centroid. If its dimension exceeds one, find a nontrivial idempotent and
split its image and complementary image, recursively. This manufactures
all primitive components in polynomial time. A third tensor has no
cross-component term, by moving an idempotent between slots. Its complete
cross-component quadratic part remains explicit.

Acquire and fit each rank-three chain component as above. Subtract the
block-diagonal compatible quadratic bases and compute the actual global
defect rank $r$. The same $2^r$ character expansion leaves a product of
independent exact component counts. Thus any number of acquired components
is permitted, with time $2^r\poly(L+n)$ and polynomial space.
When errors occur only between components, no local decoding-radius
condition is needed: all within-component bases are already exact.
The primitive split and the global coupling are both acquired from raw ANF.

\subsection{Gate descriptions with many intermediate states}
The executed front end accepts an acyclic fan-in-two CONST/XOR/AND/NOT
program with $n$ independent bits and $S$ internal gates. It manufactures
the full ANF of every wire, checks a term budget $M$ after exact
cancellation, and requires only the final polynomial to have degree at
most three. Sparse high-degree intermediates are permitted.
AND multiplication sorts at most $M^2$ union masks and cancels parity;
no high-degree terms are silently discarded.

The circuit description has $O((n+S)\log(n+S))$ bits. A conservative
front-end cost is
\[
 O\bigl(|C|+n^2+SM^2n\log(M+2)\bigr)
\]
bit time and $O(|C|+n^2+(SM+M^2)n)$ space, followed by
\eqref{eq:cost}. The implemented default is
$M=\sum_{j=0}^3\binom nj$. Exceeding the budget rejects the input.
No internal wire becomes an independent counting variable. This is a
verified symbolic front end for accepted programs, not a claim that
arbitrary programs have small ANFs. The preserved preceding work proves
that recognizing Boolean output degree at most three for unrestricted
circuits is coNP-complete.

\section{Primary literature and the claimed contribution}
The governing OpenAI manuscript constructs replacement forms and shared
linear summaries and proves subexponential exact counting for each fixed
clause width \cite{eth}. Its sampling, graph-counting, constant-locality
inversion and tensor consequences are part of that baseline. None is
claimed here as new. The present algorithms instead acquire a different
representation directly from an accepted raw polynomial or gate description;
they do not introduce one independent SAT variable per internal wire.
They do not improve the governing theorem for arbitrary fixed-width CNF.

Tensor centroids and rewriting tensors over their centroids are established
computational-algebra machinery, documented by Magma's centroid and
\texttt{TensorOverCentroid} interfaces \cite{magma}.
Support intersections have low-rank parity-check decoding ancestry;
Burle--Otmani also explains why naive independent-subspace probability
arguments can be unreliable \cite{burleotmani}. The global defect proof
here uses its own explicit direct-sum condition and union bound.
Its left/right annihilator has classical displacement-operator ancestry \cite{displacement}.
These components are not asserted to be newly invented.

Local cancellation and nondegenerate quadratic sums belong to the
stationary-phase tradition. Z\'u\~niga-Galindo proves effective rationality
for a semiquasihomogeneous local-zeta family using stationary phase
\cite{zuniga}; Sawin develops wild local Kloosterman analysis \cite{sawin}.
The compact-root-tree algorithm of Chakrabarti--Saxena is a substantive
comparison: it gives polynomial dependence on $p+k$ for fixed degree and
number of variables over $\mathbb Z/p^k\mathbb Z$, including singular
roots \cite{chakrabarti}. Its stated domain is not
$\F_{2^d}[t]/t^e$ with growing $d$, and a direct replacement of $p$ by
$q=2^d$ would leave an exponential dependence on $d$. A congruence-zero
counter also needs an additional argument to evaluate a specified primitive
character. These differences are qualifications, not claims of a lower
bound for adapting that work.

Lauder--Wan already provides the uniform finite-field terminal
\cite{lauderwan}. The candidate contribution of the present construction
is the verified chain of interfaces: raw centroid acquisition, exact
homogeneous fitting modulo a manufactured quotient, affine normalization,
constant-size projective critical-point acquisition by binary equations,
quadratic radial closure with arbitrary compatible noise, and exact
reconstruction whose costs are measured against original input bits.
The proof and executions establish those interfaces. Historical priority
has not been established, and no practical speed advantage is claimed.

The N1077 planar trivalent tensor compiler is retained as relevant Astra
context \cite{n1077}. It acquires compatible parity gauges on its restricted
class and preserves observables by exact contraction. No reduction from
the general ring phases here to that planar admissible class has been
proved, so its FKT/CRT backend is not used as a generic contraction shortcut.
The supplied task-directed optimization proof and source receipt remain
preserved byte-for-byte. The older exact/noise-aware Gram reconstruction
source was not recovered through the selected retrieval routes; the
preservation note records that gap rather than substituting the newer
proof. This work does not assume their network model.

\section{Executed evidence and implementation boundaries}
All numerical statements in this section come from the included exact
programs and JSON outputs. Raw acquisition means that the counter received
only the original Boolean ANF (or the actual gate program). Planted
coordinates were used solely by independent identity/count oracles.
No benchmark here establishes a practical speed advantage.

\begin{center}
\small
\begin{tabular}{lrrrrl}
\toprule
Executed raw family & $n$ & $d$ & $e$ & noise rank & reconstruction\\
\midrule
Smooth quotient &72&12&2&36& no field or Walsh leaves\\
Three singular rays &72&12&2&36& six radial Gaussians\\
Deep smooth quotient &108&3&12&36& origin reductions only\\
Adaptive extra digit &36&3&4&18& nine-bit residual\\
Adaptive extra digit &63&3&7&26& nine-bit residual\\
Global field defect &72&12&2&8& 256 streamed Walsh leaves\\
Global nilpotent defect &108&3&12&10& 1024 streamed Walsh leaves\\
Global field defect &144&12&4&10& 1024 streamed Walsh leaves\\
\bottomrule
\end{tabular}
\end{center}

The raw three-ray example gives
\[
 Z=4260607557632,\qquad
 N_0=2361183243565126385664.
\]
Its actual gate description has $72$ independent inputs and $4888$
internal gates ($882$ AND, $4004$ XOR and two constants). Exact gate
compilation reached $31425$ intermediate ANF terms and $31234$ output
terms. The complete output coefficient set matches the independent raw
fixture; another $1024$ original-input evaluations match. Gate states
are never added to the independent-variable count.

A second actual program has $45$ original inputs, $6379$ internal gates,
$7722$ maximum intermediate terms and $7517$ final terms. It gives
$Z=-134217728$ and $N_0=17592118935552$. The older $42$-input,
$2369$-gate split-ring program remains executable and agrees with its
previous exact result. These programs include the affine disguise and
ring arithmetic as real gates; they are not merely ANFs relabeled as
circuits.

The smooth raw suite has seven cases through $108$ bits. The separate
raw homogeneous-quotient suite has seven cases through $72$ bits,
including nonsplit nodal and three-ray reductions. Every case has $1024$
additional representation identities; its small cases also have complete
truth-table counts. The new quotient fitter was separately run on all
seven acquired smooth representations. Minimum-height acquisition passed
three raw cases, including the full-height $18$-bit case with a complete
truth table and the $36$- and $63$-bit nine-bit residuals.

The main coefficient-level weighted-ray suite has $88$ cases and
$291836416$ truth assignments across those cases. An independently
organized unweighted suite has $88$ cases and $6628416$ assignments.
A separate implementation passed $44$ full-enumeration cases and three
$30$-bit examples against a low-prefix/quadratic-fiber oracle.
Further independent audits checked $32$ fixed-unit transverse sums,
$31$ radial sums, $24$ projective coverage instances, $14$ weighted
transverse sums, two rays revived by a quadratic weight, seven entire
fitting kernels, and $5148$ height/depth bounds. These tests target
individual identities as well as final totals.

The larger $1152$-bit split-product and $12288$-bit Kloosterman-derived
examples validate coefficient backends only. They are not claims of raw
acquisition at those sizes. Conversely the $72$- and $108$-bit examples
identified as raw above include centroid and coordinate acquisition.
The general finite-field point-counting backend is a proved, cited
algorithmic dependency, not an executed component of this package.
The code rejects a leaf beyond its explicit enumeration budget unless
an exact external terminal is supplied.

\clearpage
\begin{thebibliography}{12}
\bibitem{eth} OpenAI. \emph{The Exponential Time Hypothesis is False}.
Governing release, 25 October 2026, commit 7ef5c0cc5ca4b0e072b7a050fca09ab24bc6fa3a.
\href{https://github.com/openai/math/blob/7ef5c0cc5ca4b0e072b7a050fca09ab24bc6fa3a/preprints/The-Exponential-Time-Hypothesis-is-False-October-25-2026/paper.pdf}{Pinned manuscript}.
\bibitem{lauderwan} Alan G. B. Lauder and Daqing Wan.
\emph{Counting points on varieties over finite fields of small characteristic}.
Proposition 36 and Section 6.4. \url{https://arxiv.org/pdf/math/0612147}.
\bibitem{sawin} Will Sawin.
\emph{The size of wild Kloosterman sums in number fields and function fields}.
Version 2, 31 May 2023. \url{https://arxiv.org/pdf/2209.02170v2}.
\bibitem{magma} Computational Algebra Group.
\emph{Magma Handbook: Linear Invariants of Tensors}.
\url{https://magma.maths.usyd.edu.au/magma/handbook/text/672}.
\bibitem{burleotmani} \'Etienne Burle and Ayoub Otmani.
\emph{An Upper-Bound on the Decoding Failure Probability of the LRPC Decoder}.
\url{https://arxiv.org/pdf/2309.14028}.
\bibitem{displacement} Thomas Kailath, Sun-Yuan Kung and Martin Morf.
\emph{Displacement ranks of matrices and linear equations}.
Journal of Mathematical Analysis and Applications 68(2), 395--407, 1979.
\href{https://doi.org/10.1016/0022-247X(79)90124-0}{DOI: 10.1016/0022-247X(79)90124-0}.
\bibitem{zuniga} W. A. Z\'u\~niga-Galindo.
\emph{Igusa's local zeta functions of semiquasihomogeneous polynomials}.
\url{https://arxiv.org/pdf/math/9909184}.
\bibitem{chakrabarti} Sayak Chakrabarti and Nitin Saxena.
\emph{A compact tree structure for the roots of multivariates mod $p^k$
and the related Igusa's local zeta function}. Author-hosted manuscript.
\url{https://www.cse.iitk.ac.in/users/nitin/papers/IZF-n-var.pdf}.
\bibitem{shoup} Victor Shoup.
\emph{New algorithms for finding irreducible polynomials over finite fields}.
Mathematics of Computation 54, 435--447, 1990.
\url{https://www.shoup.net/papers/detirred.pdf}.
\bibitem{n1077} Astra Research Knowledge. \emph{N1077 packet}.
Commit \texttt{e1e2272165c8} (the linked packet pins the full revision).
\href{https://github.com/harjassand/astra-research-knowledge/blob/e1e2272165c8548326c31bb11817c28f871d21e9/knowledge/current/packets/N1077.json}{Pinned representation-compiler packet}.
\bibitem{luomoments} Jinquan Luo.
\emph{On Binary Cyclic Codes with Five Nonzero Weights}.
Sections 2--3. \url{https://arxiv.org/html/0904.2237v1}.
\bibitem{kasami} Jinquan Luo, Yuansheng Tang and Hongyu Wang.
\emph{Cyclic Codes and Sequences: the Generalized Kasami Case}.
Equations (3)--(4) and Lemma 3. \url{https://arxiv.org/pdf/0902.4510}.
\end{thebibliography}

\end{document}
